% ==========================================================
%  MASTER SUMMARY TEMPLATE — AI / ML / MATHEMATICS
%  Design Philosophy:
%    1. Readable running text is the default. Prose over boxes.
%    2. Boxes are used sparingly as visual landmarks for skimming.
%    3. Visual hierarchy: chapter > section > subsection > prose.
%    4. Every box is breakable — no orphaned whitespace at page ends.
%    5. A single consistent accent color ties all structure together.
% ==========================================================

\documentclass[11pt, oneside]{book}


% ══════════════════════════════════════════════════════════
%  PAGE GEOMETRY
% ══════════════════════════════════════════════════════════
\usepackage[
  a4paper,
  top    = 2.6cm,
  bottom = 2.1cm,
  left   = 2.5cm,
  right  = 2.5cm,
  headheight = 15pt
]{geometry}


% ══════════════════════════════════════════════════════════
%  TYPOGRAPHY
% ══════════════════════════════════════════════════════════
\usepackage[T1]{fontenc}
\usepackage[utf8]{inputenc}
\usepackage{lmodern}         % Latin Modern: improved Computer Modern
\usepackage{microtype}       % Microtypographic refinements (better justification,
                             %   fewer overfull hboxes — invaluable for math text)

% If inconsolata is not installed, simply remove this line.
% It provides a crisper monospace font for code listings.
\usepackage{inconsolata}

\setlength{\parindent}{0pt}
\setlength{\parskip}{0.62em}


% ══════════════════════════════════════════════════════════
%  MATHEMATICS
% ══════════════════════════════════════════════════════════
\usepackage{amsmath, amssymb, amsthm, mathtools}
\usepackage{bm}        % Bold math symbols: \bm{\theta}
\usepackage{bbm}       % Blackboard bold for indicator: \mathbbm{1}
\usepackage{mathrsfs}  % Script letters: \mathscr{F}, \mathscr{H}
\usepackage{nicefrac}  % Compact inline fractions: \nicefrac{1}{2}
\usepackage{cancel}    % Strike-through in math: \cancel{x}


% ══════════════════════════════════════════════════════════
%  SCIENTIFIC NOTATION
% ══════════════════════════════════════════════════════════
% siunitx — consistent typesetting of numbers and SI units.
%   Usage: \SI{1.5}{\micro\meter}  \SI{50}{\kilo\hertz}  \num{6.022e23}
\usepackage{siunitx}
\sisetup{per-mode = symbol}   % render "per unit" as e.g. m/s rather than m s^{-1}

% mhchem — chemical formulas and reaction equations.
%   Usage: \ce{H2O}   \ce{ATP -> ADP + P_i}   \ce{^{14}C}
\usepackage[version=4]{mhchem}


% ══════════════════════════════════════════════════════════
%  FIGURES, TABLES, LISTS
% ══════════════════════════════════════════════════════════
\usepackage{graphicx}
\usepackage{booktabs}         % \toprule, \midrule, \bottomrule
\usepackage{array}            % Extended column types in tabular
\usepackage{float}            % [H] float placement
\usepackage{enumitem}

\setlist[itemize]{topsep=3pt, itemsep=2pt, leftmargin=2.2em}
\setlist[enumerate]{topsep=3pt, itemsep=2pt, leftmargin=2.4em}

\usepackage{needspace}        % keep headings together with their first lines

% ── Course addition: self-drawn figures (TikZ / pgfplots) ──
\usepackage{tikz}
\usetikzlibrary{arrows.meta, positioning, shapes.geometric, calc, fit,
                backgrounds, decorations.pathreplacing}
\usepackage{pgfplots}
\pgfplotsset{compat=1.18}
\usepgfplotslibrary{groupplots, fillbetween}

% Float pages: stack floats from the top instead of spreading them out.
\makeatletter
\setlength{\@fptop}{0pt}
\setlength{\@fpsep}{14pt plus 4pt}
\setlength{\@fpbot}{0pt plus 1fil}
\makeatother


% ══════════════════════════════════════════════════════════
%  ALGORITHMS / PSEUDOCODE
% ══════════════════════════════════════════════════════════
% Remove this block if you do not need pseudocode.
\usepackage[ruled, vlined, linesnumbered]{algorithm2e}
\SetAlFnt{\small}
\SetAlCapFnt{\small\bfseries}
\SetCommentSty{textit}


% ══════════════════════════════════════════════════════════
%  COLORS  (a refined, academic palette)
% ══════════════════════════════════════════════════════════
\usepackage{xcolor}

% Primary structural accent — deep navy used for headings, rules, ToC
\definecolor{accent}{HTML}{1B3A5C}

% Box accent colors: frame and background, designed to be legible but
% unobtrusive so running text remains the visual focus.
\definecolor{def@frame}{HTML}{1A5276}   \definecolor{def@bg}{HTML}{EBF5FB}
\definecolor{thm@frame}{HTML}{922B21}   \definecolor{thm@bg}{HTML}{FDEDEC}
\definecolor{ex@frame} {HTML}{1E8449}   \definecolor{ex@bg} {HTML}{EAFAF1}
\definecolor{int@frame}{HTML}{9A6A00}   \definecolor{int@bg}{HTML}{FEF9E7}
\definecolor{note@frame}{HTML}{5B2C6F}  \definecolor{note@bg}{HTML}{F5EEF8}


% ══════════════════════════════════════════════════════════
%  CODE LISTINGS
% ══════════════════════════════════════════════════════════
\usepackage{listings}
\lstset{
  basicstyle       = \ttfamily\small,
  breaklines       = true,
  numbers          = left,
  numberstyle      = \tiny\color{gray!70},
  frame            = tb,
  framerule        = 0.4pt,
  rulecolor        = \color{gray!35},
  tabsize          = 2,
  showstringspaces = false,
  keywordstyle     = \color{def@frame}\bfseries,
  commentstyle     = \color{gray!55!black}\itshape,
  stringstyle      = \color{ex@frame!90!black},
  backgroundcolor  = \color{gray!5},
  xleftmargin      = 2em,
  xrightmargin     = 0.4em,
  aboveskip        = 0.8em,
  belowskip        = 0.5em
}


% ══════════════════════════════════════════════════════════
%  CUSTOM BOXES  (tcolorbox)
%
%  KEY DESIGN DECISIONS:
%  1. enhanced   — enables advanced tcolorbox features (borderline etc.)
%  2. breakable  — box may split across a page break; eliminates the ugly
%                  white gap that appears when a long box would overflow.
%  3. Left-accent bar style — a thick colored stripe on the left margin.
%                  This is legible, modern, and unobtrusive relative to
%                  a full colored frame. The box background is a very
%                  light tint, keeping the reading experience calm.
%  4. sharp corners=west — left side is flat (squared), right is rounded.
%                  This mirrors the visual logic of the left accent bar.
%  5. parbox=false — allows display math to stretch properly inside boxes.
% ══════════════════════════════════════════════════════════
\usepackage[most]{tcolorbox}

\tcbset{
  enhanced,
  breakable,
  arc             = 2.5mm,
  boxrule         = 0.5pt,
  left            = 4.5mm,
  right           = 4mm,
  top             = 2.5mm,
  bottom          = 2.5mm,
  fonttitle       = \bfseries\small,
  sharp corners   = west,
  parbox          = false
}

\newtcolorbox{defbox}[1][]{
  title  = {Definition},
  colback = def@bg, colframe = def@frame,
  borderline west = {3.5pt}{0pt}{def@frame},
  #1
}

\newtcolorbox{thmbox}[1][]{
  title  = {Theorem / Key Formula},
  colback = thm@bg, colframe = thm@frame,
  borderline west = {3.5pt}{0pt}{thm@frame},
  #1
}

\newtcolorbox{exbox}[1][]{
  title  = {Example},
  colback = ex@bg, colframe = ex@frame,
  borderline west = {3.5pt}{0pt}{ex@frame},
  #1
}

\newtcolorbox{intbox}[1][]{
  title  = {Intuition \& Mental Model},
  colback = int@bg, colframe = int@frame,
  borderline west = {3.5pt}{0pt}{int@frame},
  #1
}

\newtcolorbox{notebox}[1][]{
  title  = {Note},
  colback = note@bg, colframe = note@frame,
  borderline west = {3.5pt}{0pt}{note@frame},
  #1
}


% ══════════════════════════════════════════════════════════
%  TABLE OF CONTENTS  (tocloft)
%
%  The [titles] option prevents tocloft from conflicting
%  with titlesec's redefinition of chapter/section titles.
% ══════════════════════════════════════════════════════════
\usepackage[titles]{tocloft}

% Chapter entries: accented, bold, with dots leading to page number
\renewcommand{\cftchapfont}     {\bfseries\color{accent}}
\renewcommand{\cftchappagefont} {\bfseries\color{accent}}
\renewcommand{\cftchapleader}   {\bfseries\color{accent!40}\cftdotfill{\cftdotsep}}
\setlength{\cftbeforechapskip} {7pt}

% Section entries: standard weight, small size
\renewcommand{\cftsecfont}      {\small}
\renewcommand{\cftsecpagefont}  {\small}
\setlength{\cftsecindent}      {1.8em}
\setlength{\cftbeforesecskip}  {1.5pt}

% Subsection entries: muted gray to visually recede
\renewcommand{\cftsubsecfont}      {\small\color{gray!65!black}}
\renewcommand{\cftsubsecpagefont}  {\small\color{gray!65!black}}
\setlength{\cftsubsecindent}      {3.4em}
\setlength{\cftbeforesubsecskip}  {0.5pt}


% ══════════════════════════════════════════════════════════
%  CHAPTER & SECTION HEADINGS  (titlesec)
%
%  Chapter: large faded chapter number floated right, title
%           left-aligned beneath it, then a double rule.
%           The oversized ghost number is a quiet nod to
%           modern textbook design without being gimmicky.
%
%  Section: accent-colored, with a thin rule beneath.
%           The rule draws the eye without shouting.
%
%  Subsection: slightly smaller accent color, no rule — the
%           rule at section level is already enough structure.
% ══════════════════════════════════════════════════════════
\usepackage{titlesec}

\titleformat{\chapter}[display]
  {\normalfont\bfseries\color{accent}\huge}
  {\hfill\fontsize{68}{68}\selectfont\color{accent!13}\thechapter}
  {-10pt}
  {}
  [{\vspace{4pt}%
    {\color{accent}\rule{\linewidth}{1.6pt}}\\[-2pt]%
    {\color{accent!35}\rule{\linewidth}{0.5pt}}}]
\titlespacing*{\chapter}{0pt}{-28pt}{22pt}

\titleformat{\section}
  {\normalfont\large\bfseries\color{accent}}
  {\thesection}
  {0.8em}
  {}
  [{\vspace{1pt}\color{accent!45}\titlerule[0.45pt]}]
\titlespacing*{\section}{0pt}{3.8ex plus 1ex minus .2ex}{2.2ex plus .2ex}

\titleformat{\subsection}
  {\normalfont\normalsize\bfseries\color{accent!80!black}}
  {\thesubsection}
  {0.7em}
  {}
\titlespacing*{\subsection}{0pt}{2.6ex plus 1ex minus .2ex}{1.4ex plus .2ex}

% Run-in paragraph label — useful for named sub-items in a derivation.
\titleformat{\paragraph}[runin]
  {\normalfont\small\bfseries}
  {}
  {0em}
  {}[.\enspace]
\titlespacing{\paragraph}{0pt}{1.2ex plus .5ex minus .2ex}{0.5em}


% ══════════════════════════════════════════════════════════
%  HEADER & FOOTER
% ══════════════════════════════════════════════════════════
\usepackage{fancyhdr}
\pagestyle{fancy}
\fancyhf{}
\fancyhead[L]{\small\color{gray!70}\nouppercase{\leftmark}}
\fancyhead[R]{\small\color{gray!70}\thepage}
\renewcommand{\headrulewidth}{0.4pt}

% Chapter-opening pages use LaTeX's 'plain' style — redefine it here
% to match the rest of the document (page number at foot, no header line).
\fancypagestyle{plain}{%
  \fancyhf{}
  \fancyfoot[C]{\small\color{gray!70}\thepage}
  \renewcommand{\headrulewidth}{0pt}%
}


% ══════════════════════════════════════════════════════════
%  LINKS & REFERENCES
% ══════════════════════════════════════════════════════════
\usepackage[hidelinks]{hyperref}
\usepackage[nameinlink, capitalize]{cleveref}


% ══════════════════════════════════════════════════════════
%  DOCUMENT METADATA  (fill in for each new subject)
% ══════════════════════════════════════════════════════════
\newcommand{\coursetitle}{Deep Learning: Basic Techniques}
\newcommand{\documenttype}{Study Notes -- Working Draft}
\newcommand{\authorname}{Stefan Eder}
\newcommand{\basedon}{Based on the lecture \emph{Deep Learning: Basic Techniques} (G.~Klambauer, JKU Linz), slides of WS 2025/26}


% ══════════════════════════════════════════════════════════
%  MATHEMATICAL MACROS
%  Organized by domain. Add your own at the end of each group.
% ══════════════════════════════════════════════════════════

% ── Number Sets & Spaces ──────────────────────────────────
\newcommand{\R}{\mathbb{R}}
\newcommand{\N}{\mathbb{N}}
\newcommand{\Z}{\mathbb{Z}}
\newcommand{\C}{\mathbb{C}}
\newcommand{\Q}{\mathbb{Q}}

% ── Probability & Statistics ──────────────────────────────
\newcommand{\E}{\mathbb{E}}              % Expectation
\newcommand{\Prob}{\mathbb{P}}           % Probability measure
\newcommand{\Var}{\mathrm{Var}}
\newcommand{\Cov}{\mathrm{Cov}}
\newcommand{\Corr}{\mathrm{Corr}}
\newcommand{\iid}{\overset{\mathrm{iid}}{\sim}}
% Indicator function. Usage: \ind_{\{x > 0\}}
\newcommand{\ind}{\mathbbm{1}}

% ── Information Theory ────────────────────────────────────
% KL divergence. Usage: \KL{q}{p}  →  D_KL(q ‖ p)
\newcommand{\KL}[2]{D_{\mathrm{KL}}\!\left(#1 \,\|\, #2\right)}
\newcommand{\Ent}{\mathrm{H}}            % Entropy:         \Ent(X)
\newcommand{\MI}{\mathrm{I}}             % Mutual information: \MI(X;Y)

% ── Calculus & Optimization ───────────────────────────────
\newcommand{\argmax}{\operatorname*{arg\,max}}
\newcommand{\argmin}{\operatorname*{arg\,min}}
\newcommand{\grad}{\nabla}               % Gradient
\newcommand{\hess}{\nabla^2}             % Hessian
\newcommand{\lapl}{\Delta}               % Laplacian
% Differential (upright d). Usage: \int f(x) \dd x
\newcommand{\dd}{\,\mathrm{d}}
% Partial derivatives. Usage: \pd{f}{x}  or  \pdd{f}{x}
\newcommand{\pd}[2]{\frac{\partial #1}{\partial #2}}
\newcommand{\pdd}[2]{\frac{\partial^2 #1}{\partial #2^2}}

% ── Norms, Abs, Inner Products ────────────────────────────
\newcommand{\norm}[1]{\left\lVert #1 \right\rVert}
\newcommand{\normsq}[1]{\left\lVert #1 \right\rVert^2}
\newcommand{\abs}[1]{\left\lvert #1 \right\rvert}
\newcommand{\inner}[2]{\left\langle #1,\, #2 \right\rangle}

% ── Linear Algebra ────────────────────────────────────────
\newcommand{\T}{^{\mathsf{T}}}           % Transpose:        \mathbf{A}\T
\newcommand{\inv}{^{-1}}                 % Inverse:          A\inv
\newcommand{\pinv}{^{\dagger}}           % Pseudo-inverse:   A\pinv
\newcommand{\tr}{\operatorname{tr}}
\newcommand{\rank}{\operatorname{rank}}
\newcommand{\diag}{\operatorname{diag}}
\newcommand{\spanop}{\operatorname{span}}
\newcommand{\eye}{\mathbf{I}}            % Identity matrix
\newcommand{\zeros}{\mathbf{0}}
\newcommand{\ones}{\mathbf{1}}

% ── ML / AI Core Notation ────────────────────────────────
\newcommand{\loss}{\mathcal{L}}          % Loss:          \loss(\params)
\newcommand{\reg}{\mathcal{R}}           % Regularizer:   \reg(\params)
\newcommand{\dataset}{\mathcal{D}}       % Dataset:       \dataset
\newcommand{\params}{\boldsymbol{\theta}}   % Parameters
\newcommand{\weights}{\mathbf{w}}
\newcommand{\xvec}{\mathbf{x}}
\newcommand{\yvec}{\mathbf{y}}
\newcommand{\zvec}{\mathbf{z}}
\newcommand{\hvec}{\mathbf{h}}           % Hidden state
\newcommand{\feature}{\boldsymbol{\phi}} % Feature map

% Activation functions
\DeclareMathOperator{\softmax}{softmax}
\DeclareMathOperator{\relu}{ReLU}
\DeclareMathOperator{\sign}{sign}
% sigmoid: just use \sigma (standard Greek letter)

% ── Complexity ────────────────────────────────────────────
\newcommand{\bigO}[1]{\mathcal{O}\!\left(#1\right)}
\newcommand{\bigOm}[1]{\Omega\!\left(#1\right)}
\newcommand{\bigTh}[1]{\Theta\!\left(#1\right)}

% ── Common Distributions ──────────────────────────────────
\newcommand{\Normal}{\mathcal{N}}
\newcommand{\Uniform}{\mathcal{U}}
\newcommand{\Bernoulli}{\mathrm{Ber}}
\newcommand{\Categorical}{\mathrm{Cat}}
\newcommand{\Dirichlet}{\mathrm{Dir}}
\newcommand{\Binomial}{\mathrm{Bin}}

% ── Miscellaneous ─────────────────────────────────────────
\newcommand{\given}{\,\vert\,}           % Conditional: \Prob(A \given B)
\newcommand{\st}{\;\text{s.t.}\;}        % "subject to"
\newcommand{\eqdef}{\coloneqq}           % Definitional equality  :=
\newcommand{\eqq}{\overset{?}{=}}        % "Does this equal?"

% ══════════════════════════════════════════════════════════
%  COURSE-SPECIFIC MACROS  (Deep Learning: Basic Techniques)
%  Notation follows the lecture slides: bold italic vectors and
%  matrices, sample index as superscript (x^n), layer index in
%  square brackets (W^{[l]}), weight w_{ij} from unit j to unit i.
% ══════════════════════════════════════════════════════════
\newcommand{\bx}{\bm{x}}        \newcommand{\by}{\bm{y}}
\newcommand{\bw}{\bm{w}}        \newcommand{\bW}{\bm{W}}
\newcommand{\ba}{\bm{a}}        \newcommand{\bs}{\bm{s}}
\newcommand{\bb}{\bm{b}}        \newcommand{\bh}{\bm{h}}
\newcommand{\bg}{\bm{g}}        \newcommand{\be}{\bm{e}}
\newcommand{\bX}{\bm{X}}        \newcommand{\bY}{\bm{Y}}
\newcommand{\bI}{\bm{I}}        \newcommand{\bJ}{\bm{J}}
\newcommand{\bH}{\bm{H}}        \newcommand{\bR}{\bm{R}}
\newcommand{\bdelta}{\bm{\delta}}
\newcommand{\byhat}{\hat{\bm{y}}}
\newcommand{\what}{\hat{\bm{w}}}
\newcommand{\lay}[1]{^{[#1]}}            % layer index:  \bW\lay{l}
\newcommand{\Remp}{R_{\mathrm{emp}}}     % empirical error
\newcommand{\Rgen}{R_{\mathrm{gen}}}     % generalization error (Ch. 5)
\newcommand{\dsig}{\sigma'}              % derivative of the sigmoid
\DeclareMathOperator{\tanhop}{tanh}

% ── Figure colours (aliases of the template palette) ──────
\colorlet{figblue}{def@frame}
\colorlet{figred}{thm@frame}
\colorlet{figgreen}{ex@frame}
\colorlet{figamber}{int@frame}
\colorlet{figpurple}{note@frame}

% ── TikZ styles for network diagrams ──────────────────────
\tikzset{
  unit/.style    = {circle, draw=accent, thick, fill=accent!6,
                    minimum size=8mm, inner sep=0pt, font=\small},
  inunit/.style  = {unit, fill=gray!12, draw=gray!70!black},
  outunit/.style = {unit, fill=accent!18},
  conn/.style    = {-{Stealth[length=2.2mm]}, accent!80!black, thin},
  bconn/.style   = {-{Stealth[length=2.4mm]}, figred, thick},
  lbl/.style     = {font=\scriptsize, inner sep=1pt, fill=white},
  every pin/.style = {font=\scriptsize},
}

% ── pgfplots look shared by all plots ─────────────────────
\pgfplotsset{
  dlbt/.style = {
    axis lines = middle,
    axis line style = {gray!70, -{Stealth[length=1.6mm]}},
    tick label style = {font=\scriptsize, gray!70!black},
    label style = {font=\small},
    grid = major, grid style = {gray!15},
    samples = 160,
  },
}

% ── Placeholders ──────────────────────────────────────────
% Figure placeholder (same signature as in other course notes).
% #1 height, #2 concept, #3 what the picture should show + search hint,
% #4 caption, #5 label.  Replace with \includegraphics{figures/...}
% once a suitable image is in figures/.
\newcommand{\figureplaceholder}[5]{%
\begin{figure}[htbp]
\centering
\begin{tcolorbox}[unbreakable,width=\linewidth,height=#1,colback=accent!3,
  colframe=accent!30,boxrule=0.5pt,arc=2.5mm,sharp corners=all,
  valign=center,halign=center,fontupper=\small]
{\bfseries\color{accent}Figure to insert: #2}\par\medskip
{\color{gray!75!black}#3}
\end{tcolorbox}
\caption{#4}\label{#5}
\end{figure}}

% Outline text for chapters that are not written yet.
\newcommand{\planned}[1]{%
  \par{\small\color{gray!70!black}\textit{Planned content:} #1}\par}

% Keep section headings together with the start of their text.
\let\originalsection\section
\renewcommand{\section}{\Needspace{7\baselineskip}\originalsection}
\let\originalsubsection\subsection
\renewcommand{\subsection}{\Needspace{6\baselineskip}\originalsubsection}


% ==========================================================
%  BEGIN DOCUMENT
% ==========================================================
\begin{document}

\frontmatter  % Roman-numeral page numbers for front matter


% ──────────────────────────────────────────────────────────
%  TITLE PAGE
% ──────────────────────────────────────────────────────────
\thispagestyle{empty}
\vspace*{1.8cm}

\begin{center}
  {\color{accent}\rule{\linewidth}{1.8pt}}
  \vspace{0.65cm}

  {\fontsize{28}{34}\selectfont\bfseries\color{accent}\coursetitle\par}
  \vspace{0.25cm}
  {\large\color{gray!80!black}\documenttype\par}

  \vspace{0.55cm}
  {\color{accent!35}\rule{\linewidth}{0.5pt}}
  \vspace{0.55cm}

  {\large\authorname\par}
  \vspace{0.2cm}
  {\small\color{gray!70}\basedon\par}
\end{center}

\vspace{1.6cm}

\begin{center}
\begin{minipage}{0.84\textwidth}
  \small\color{gray!80!black}
  These notes walk through the basic techniques of deep learning the way the
  lecture builds them up: from a single linear neuron, via the perceptron,
  logistic and softmax regression, to multi-layer and deep networks trained
  with backpropagation. Every idea comes in plain words first, then as a
  formula, and wherever it helps there is a small example with actual numbers.
  All worked numbers were re-computed by script, and the backpropagation
  example was checked against finite differences.
  Use the notes as a study aid alongside the lecture slides, not as a
  replacement for them.

  \medskip
  \textbf{Working draft:} Chapters 1, 3 and 4 are written out. Chapter 2 and
  Chapters 5--10 are still outlines and will be filled in once their slides
  are available.

  \medskip
  If you spot an error or a gap, please open an issue or pull request at
  \textbf{github.com/steveTgit/AI\_Master} --- corrections are always welcome.

  \medskip
  \color{gray!60!black}
  Good study materials should be shared, not gatekept.
  If this helped you, take it as a reason to help the next colleague you see struggling with something.
\end{minipage}
\end{center}

\vfill
{\centering\small\color{gray!65} \today\par}
\clearpage


% ──────────────────────────────────────────────────────────
%  TABLE OF CONTENTS
% ──────────────────────────────────────────────────────────
\setcounter{tocdepth}{2}  % Show: chapters, sections, subsections only

\begingroup
  \let\cleardoublepage\relax
  \let\clearpage\relax
  \tableofcontents
\endgroup

\mainmatter  % Switch to Arabic numerals


% ==========================================================
%  CHAPTER 1 — A BRIEF HISTORY OF NEURAL NETWORKS
%  Source: "Deep Learning and Neural Networks I: Intro & History of NN"
% ==========================================================
\chapter{A Brief History of Neural Networks and Deep Learning}
\label{ch:history}

The first lecture does two things. It sets up the course (what it covers,
how it is examined, what to read), and it tells the story of how neural
networks went from a 1940s model of a nerve cell to the technology behind
almost every recent AI headline. The history is worth knowing for its own
sake, but it also previews the technical content: nearly every milestone in
this chapter --- the perceptron, the XOR problem, backpropagation, the
vanishing gradient, ReLUs, convolutions --- comes back later in these notes
as a formula.

\section{About the Course}

\emph{Deep Learning: Basic Techniques} (some slide decks use the title
\emph{Deep Learning and Neural Networks I}) is a first-semester basic course
of the AI master at JKU Linz. It consists of a lecture (VO, 2~hours per week,
3~ECTS), examined by a written on-site exam at the end of the semester, and
an exercise (UE, 1~hour per week, 1.5~ECTS) with five to six roughly
bi-weekly assignments. The lecture notes on KUSSS/Moodle contain a reference
list; the main textbooks recommended in the first lecture are:
\begin{itemize}
  \item C.\,M.~Bishop \& H.~Bishop (2023), \emph{Deep Learning: Foundations
    and Concepts}, Springer;
  \item C.\,M.~Bishop (2006), \emph{Pattern Recognition and Machine
    Learning}, Springer;
  \item I.~Goodfellow, Y.~Bengio \& A.~Courville (2016), \emph{Deep
    Learning}, MIT Press;
  \item M.\,P.~Deisenroth, A.\,A.~Faisal \& C.\,S.~Ong (2020),
    \emph{Mathematics for Machine Learning}, Cambridge University Press;
  \item A.~Zhang, Z.\,C.~Lipton, M.~Li \& A.\,J.~Smola, \emph{Dive into Deep
    Learning} (freely available online).
\end{itemize}

\Cref{tab:roadmap} lists the chapters of the course. These notes follow the
same numbering, so ``Chapter~4'' on the slides is \cref{ch:basic} here, and
the section numbers 4.1--4.10 match as well.

\begin{table}[htbp]
\centering
\small
\begin{tabular}{@{}rll@{}}
\toprule
\textbf{Ch.} & \textbf{Topic} & \textbf{Status in these notes} \\
\midrule
1  & A brief history of neural networks          & written \\
2  & Biological neural networks                  & outline \\
3  & Notation                                    & written (collected from Ch.~4) \\
4  & Basic neural network architectures (4.1--4.10) & written \\
5  & Generalization error, test set and cross-validation & outline \\
6  & Convolutional neural networks               & outline \\
7  & Learning methods: basic algorithms          & outline \\
8  & Regularization                              & outline \\
9  & Initialization                              & outline \\
10 & Normalization                               & outline \\
\bottomrule
\end{tabular}
\caption{Roadmap of the course as announced in the first lecture.
Chapter~4 is by far the largest: it goes from the single linear neuron
(4.1) to deep feed-forward networks, backpropagation and the practical
``tricks of the trade'' (4.10).}
\label{tab:roadmap}
\end{table}

And one piece of advice from the first lecture that is not about maths:
connect with your colleagues. Most of this material is much easier to learn
in a group than alone.

\section{Machine Learning and Deep Learning}

The lecture opens with three nested circles: \emph{intelligent machines and
programs} (artificial intelligence), inside it the \emph{ability to learn}
(machine learning), and inside that \emph{neural networks} (deep learning).
The claim behind the picture is that the current success of AI is, to a
large extent, the success of deep learning.

\begin{defbox}[title={Machine learning and deep learning}]
\textbf{Machine learning} is ``the science of getting computers to act
without being explicitly programmed'' (Andrew Ng): the system infers
knowledge from data and has to \emph{generalize} to data it has not seen,
for example in supervised learning, where it learns to classify from
labelled examples.

\textbf{Deep learning} is machine learning with (deep) neural networks.
Instead of relying on hand-crafted features, deep networks \emph{learn
representations} of the data themselves, in several levels: lower levels
code simple patterns, higher levels code increasingly abstract concepts that
help to solve the task.
\end{defbox}

\begin{intbox}
The standard picture is a face-recognition network. The first layer learns
something like edge and blob detectors, the next layers combine them into
parts (eyes, noses, mouths), and the top layers respond to whole faces.
Nobody programmed ``look for an eye'' --- the hierarchy emerges because it is
useful for minimizing the training loss. That is what ``learning increasingly
abstract features'' means, and it is the main difference to classical machine
learning, where a human designs the features and the model only learns the
last step.
\end{intbox}

\section{A Brief History}
\label{sec:history}

The history in the lecture follows J.~Schmidhuber's review \emph{Deep
learning in neural networks: an overview} (Neural Networks, 2015), and the
lecture is explicit about a caveat: all history is biased, incomplete and
inaccurate --- and the history of neural networks is interwoven with, but
not the same as, the history of AI. \Cref{fig:timeline} gives the overview;
the subsections fill in the details.

\begin{figure}[htbp]
\centering
\begin{tikzpicture}[x=1.7mm, y=1cm,
    ev/.style = {font=\tiny, align=center, text width=15mm, inner sep=1pt},
    tk/.style = {gray!55, thin}]
  % winters
  \fill[gray!11] (30,-2.75) rectangle (40,2.85);
  \fill[gray!11] (50,-2.75) rectangle (65,2.85);
  \node[font=\scriptsize\itshape, gray!65!black, anchor=south] at (35,-2.75)
    {1st winter};
  \node[font=\scriptsize\itshape, gray!65!black, anchor=south] at (59.5,-2.75)
    {2nd winter};
  % axis
  \draw[thick, accent, -{Stealth[length=2.4mm]}] (0,0) -- (86,0);
  % helper: \ev{year}{level}{text}
  \newcommand{\ev}[3]{%
    \pgfmathsetmacro{\xx}{#1-1940}%
    \draw[tk] (\xx,0) -- (\xx,#2);%
    \ifdim #2pt>0pt
      \node[ev, anchor=south] at (\xx,#2) {\textbf{#1}\\#3};
    \else
      \node[ev, anchor=north] at (\xx,#2) {\textbf{#1}\\#3};
    \fi}
  % above: techniques
  \ev{1943}{0.45}{McCulloch--Pitts neuron}
  \ev{1957}{0.45}{Rosenblatt's perceptron}
  \ev{1969}{1.25}{Minsky \& Papert: XOR}
  \ev{1986}{0.45}{Backprop popularized}
  \ev{1991}{1.25}{Vanishing gradients}
  \ev{1997}{2.05}{LSTM}
  \ev{2006}{2.05}{Stacked pre-training}
  \ev{2012}{0.45}{AlexNet wins ImageNet}
  \ev{2017}{1.45}{Transformer}
  \ev{2022}{2.05}{Large language models}
  % below: more milestones
  \ev{1949}{-0.45}{Hebbian learning}
  \ev{1965}{-0.45}{GMDH: first deep nets}
  \ev{1974}{-1.25}{Backprop for NNs (Werbos)}
  \ev{1982}{-0.45}{Hopfield networks}
  \ev{1989}{-1.25}{CNNs (LeCun); universal approximation}
  \ev{2009}{-0.45}{ImageNet data set}
  \ev{2016}{-1.25}{AlphaGo}
  \ev{2021}{-0.45}{AlphaFold}
  % year ticks drawn last, so they sit on top of the leader lines
  \foreach \y in {1940,1950,...,2020} {
    \pgfmathsetmacro{\xx}{\y-1940}
    \draw[accent] (\xx,-0.08) -- (\xx,0.08);
    \node[font=\tiny, accent, fill=white, inner sep=0.6pt] at (\xx,-0.22) {\y};
  }
  \foreach \y in {1943,1949,1957,1965,1969,1974,1982,1986,1989,1991,1997,2006,2009,2012,2016,2017,2021,2022}
    \fill[accent] (\y-1940,0) circle (0.45mm);
\end{tikzpicture}
\caption{Milestones of neural-network history as told in the lecture. The
shaded bands are the two ``AI winters'' for neural networks (1970--$\sim$1980
and 1990--$\sim$2005), periods in which funding and interest dried up.
Notice that backpropagation and LSTM were both invented \emph{during} a
winter.}
\label{fig:timeline}
\end{figure}

\subsection{Before 1940: Least Squares, Gradient Descent and the Chain Rule}

Three ingredients of modern deep learning are centuries old:
\begin{itemize}
  \item \textbf{Linear regression and least squares} (Legendre 1805, Gauss
    1809). A linear regression model is exactly a single linear
    ``neuron'' (\cref{sec:linear-neuron}).
  \item \textbf{Gradient descent} (Cauchy 1847, Hadamard 1908): improve a
    guess by stepping against the gradient.
  \item \textbf{The chain rule} (Leibniz 1676, later Euler): the only
    mathematical tool behind backpropagation (\cref{sec:backprop}).
\end{itemize}
Put together, ``a linear model fitted by gradient descent on a squared error,
with gradients computed by the chain rule'' already describes a good part of
\cref{ch:basic}.

\subsection{1943--1969: The First Neurons, Hebb and the Perceptron}

In 1943, \textbf{McCulloch and Pitts} published \emph{A logical calculus of
the ideas immanent in nervous activity}, the first computational model of a
neuron. Their neuron is a basic logic unit: it takes inputs, forms the
weighted sum (dot product of inputs and weights), applies an activation
function (a threshold) and ``fires'' an output. There is no learning --- the
weights are set by hand.

In 1949, \textbf{Donald O.~Hebb} proposed a hypothesis about how biological
networks learn: the connection between two neurons that fire together is
strengthened (``cells that fire together wire together''). This
\emph{Hebbian learning} reappears in the Hopfield networks at the end of
\cref{ch:basic}.

In 1957, \textbf{Frank Rosenblatt} built the \emph{perceptron}: an electronic
device with adaptable connections, 400 photocells as input neurons, used for
image recognition --- and, crucially, with the ability to \emph{learn} its
weights from examples. Its learning rule is the subject of
\cref{sec:perceptron}.

Multi-layer networks are older than often assumed: \textbf{Oleksiy
Ivakhnenko}'s \emph{Group Method of Data Handling} (GMDH, 1965 and 1971)
already trained networks with several layers.

In 1969, \textbf{Marvin Minsky and Seymour Papert} published the book
\emph{Perceptrons}. Among other things, it showed that a single-layer
perceptron cannot solve the simple XOR problem (\cref{sec:xor}). The book
contributed to a strong reduction of research on neural networks --- the
\textbf{first AI winter} (1970--$\sim$1980). Minsky received the Turing Award
in 1969.

\subsection{1974--1989: Backpropagation and Multi-Layer Perceptrons}

The way out of the XOR dead end is a hidden layer, and the way to train a
hidden layer is backpropagation. \textbf{Paul Werbos} suggested
backpropagation for neural networks in 1974; it became widely known through
\textbf{Rumelhart, Hinton and Williams} (1986) in the \emph{Parallel
Distributed Processing} volumes, which showed how to train multi-layer
perceptrons (MLPs) with it.

In 1989, \textbf{Cybenko} and \textbf{Kurt Hornik} proved versions of the
\emph{universal function approximation theorem}: an MLP with a single hidden
layer and enough hidden units can approximate any continuous function on a
compact set arbitrarily well. So MLPs do not only solve XOR --- they can, in
principle, represent (almost) anything.

Two more architectures appear in this period. In 1982, \textbf{John Hopfield}
introduced \emph{Hopfield networks}, recurrent networks that work as
associative memories. And convolutional networks have their roots here too:
\textbf{Kunihiko Fukushima}'s \emph{Neocognitron} (1979) introduced the
architecture (but was not trained by backprop), and \textbf{Yann LeCun}
trained convolutional neural networks with backpropagation for image
recognition in 1989.

The 2024 Nobel Prize in Physics went to Hopfield (``Hopfield networks'') and
Hinton (``Boltzmann machines'') for foundational work on neural networks.

\subsection{1991--1997: The Vanishing Gradient Problem and LSTM}

\begin{defbox}[title={The vanishing gradient problem (VGP)}]
Training a neural network relies on gradients, which are propagated backward
through the layers. In deep networks (and in recurrent networks unrolled
over time) these gradients shrink --- or blow up --- exponentially with the
number of layers they pass. The deeper the network, the harder it is to
train. Discovered by \textbf{Sepp Hochreiter} in his 1991 diploma thesis;
called by the lecture ``the fundamental problem of all deep learning
methods''.
\end{defbox}

The mathematical reason is derived in \cref{sec:vgp}. For recurrent networks,
Hochreiter and Schmidhuber's \textbf{Long Short-Term Memory (LSTM)} (1997)
ameliorates the problem; for feed-forward networks, ReLUs, careful
initialization and skip connections help (\cref{sec:vgp}).

The 1990s and early 2000s were the \textbf{second AI winter} for neural
networks (1990--$\sim$2005): support-vector machines and other kernel
methods dominated machine learning, and neural networks were considered hard
to train.

\subsection{Around 2010: Re-Emergence as Deep Learning}

Deep learning took off when several developments came together:
\begin{itemize}
  \item \textbf{GPU training} of neural networks (Oh and Jung, 2004);
  \item \textbf{stacked, layer-wise training} of deep networks (Hinton et
    al., 2006);
  \item \textbf{rectified linear units}, ReLUs (Nair and Hinton, 2010);
  \item \textbf{deep CNNs} trained on GPUs (Krizhevsky et al., 2012);
  \item \textbf{large data sets}: in 2009, Fei-Fei Li launched
    \emph{ImageNet}, a database of about 14~million labelled images that
    became \emph{the} benchmark of deep learning research.
\end{itemize}
The lecture dates the starting point of the deep learning era to about
2010. In 2018, Yoshua Bengio, Geoffrey Hinton and Yann LeCun received the
Turing Award for deep learning.

\section{Why Deep Learning Took Off}

\begin{intbox}
Neural networks are not new; what changed around 2010--2012 is the
environment they run in. The lecture lists four reasons for the boom:
\textbf{fast computers}, \textbf{GPUs}, \textbf{large data sets} and
\textbf{algorithmic improvements}. A network that was hopeless to train on
a 1995 workstation with a few thousand examples becomes practical with
a GPU and a million labelled images --- and once it worked, people found
the algorithmic tricks (ReLU, better initialization, dropout, \ldots) to make
it work even better.
\end{intbox}

The ``third wave'' started with a series of records around 2012, in vision
and in speech:
\begin{itemize}
  \item \textbf{2009:} deep (LSTM-based recurrent) networks won
    international handwriting-recognition competitions (French, Arabic and
    Chinese handwriting).
  \item \textbf{2011:} deep networks won the German traffic-sign recognition
    contest with super-human performance: 0.56\,\% error, compared with
    1.16\,\% for humans and 3.86\,\% for the best non-neural competitor.
  \item \textbf{2011:} on an industry speech-recognition benchmark, deep
    networks lowered the error rate from about 24\,\% to about 16\,\% --- after
    the community had been stuck at 24\,\% for more than a decade. This is
    when IT companies became seriously interested.
  \item \textbf{2012:} deep networks won the brain-image segmentation
    challenge with super-human pixel error rate.
  \item \textbf{2012:} a deep CNN trained on GPUs (``AlexNet'') won the
    ImageNet competition with about 16\,\% error on 1.2~million images,
    against 26\,\% for the runner-up.\footnote{The record list on the slides
    dates this to 2013; the winning network was presented at the 2012
    ImageNet challenge (Krizhevsky et al., 2012), which is also the date the
    slides use elsewhere.}
  \item \textbf{2013:} deep networks won the grand challenge on mitosis
    detection, and Google researchers improved previous face-recognition
    results by 70\,\% using 16{,}000 CPU cores.
\end{itemize}

Since then the list of milestones has kept growing. Examples from the
lecture:
\begin{itemize}
  \item \textbf{2016:} AlphaGo beats a human Go champion; deep learning
    systems understand human speech well enough for everyday assistants.
  \item \textbf{2017:} deep networks classify skin cancer at the level of
    dermatologists (Esteva et al., \emph{Nature}).
  \item \textbf{2021:} AlphaFold, a Transformer-based network for protein
    structure prediction, is named scientific breakthrough of the year.
  \item \textbf{2022:} deep reinforcement learning steers the magnetic
    confinement of plasma in a tokamak fusion reactor (Degrave et al.,
    \emph{Nature}).
  \item \textbf{2022:} large language models produce human-like text --- the
    era of LLMs and of the broad uptake of deep learning methods begins.
\end{itemize}

\section{Main Types of Networks}
\label{sec:network-types}

The lecture distinguishes four main families of networks. The useful way to
remember them is by the \emph{kind of input} they are built for.

\begin{table}[htbp]
\centering
\small
\begin{tabular}{@{}p{4.6cm}p{2.6cm}p{5.9cm}@{}}
\toprule
\textbf{Network type} & \textbf{Input} & \textbf{Typical applications} \\
\midrule
Feed-forward / fully connected networks & vectors &
  standard classification problems: spam vs.\ non-spam, cancer vs.\ normal
  tissue \\
Recurrent neural networks (RNNs) & time series, sequences &
  text, speech, video; e.g.\ chatbots \\
Convolutional neural networks (CNNs) & spatial data (images) &
  image processing, image recognition \\
Transformer networks (2017) & sets (of tokens) &
  natural language processing \\
\bottomrule
\end{tabular}
\caption{The four main network families. This course covers fully connected
networks (\cref{ch:basic}) and CNNs (\cref{ch:cnn}); RNNs, LSTM and
Transformers are the topic of later courses.}
\label{tab:network-types}
\end{table}

\section{Credit, Prizes and Disputes}

The history of neural networks is not a settled story, and the lecture shows
an example of an ongoing dispute. In 2020, Geoffrey Hinton received the
Honda Prize; J\"urgen Schmidhuber publicly criticized the choice, arguing
that inventions were attributed to Hinton that were not his. He made similar
criticism of the 2018 Turing Award to Bengio, Hinton and LeCun, and of the
fact that Hopfield networks are usually credited to Hopfield alone, although
Shun-Ichi Amari had described the same kind of network in 1972 (this is why
the lecture says \emph{Amari--Hopfield} networks, \cref{sec:hopfield}).

A more neutral record is the list of IEEE Neural Networks Pioneer Award
winners, which contains most of the names in this chapter --- among them
Widrow, Kohonen and Grossberg (1991), Amari (1992), Werbos (1995), Hopfield
(1997), Hinton (1998), Fukushima (2003), Vapnik (2010), LeCun (2014),
Schmidhuber (2016), Bengio (2019) and Hochreiter (2020).

\section{Neural-Network History versus AI History}

Everything above is a history of \emph{neural networks}, not of \emph{AI}. A
history of AI would start with Alan Turing and the Turing test, and with the
1956 Dartmouth workshop (J.~McCarthy, C.~Shannon, R.~Solomonoff, T.~More,
N.~Rochester, O.~Selfridge, A.~Samuel and others), where the term
``artificial intelligence'' was coined. It would talk about symbolic
reasoning, expert systems, and IBM's Deep Blue beating Kasparov at chess in
1997. Broadly, AI went through three waves --- the birth of AI in the 1950s
and 1960s, expert systems and knowledge engineering in the 1970s and 1980s,
and machine learning with deep networks since about 2010 --- separated by two
winters.

\begin{notebox}[title={Two different sets of ``AI winter'' dates}]
The general AI history places the winters roughly in the mid-1970s and in
the late 1980s to early 1990s, triggered by limited computing power and by
the failure of expert systems to scale. The neural-network winters on the
slides (1970--$\sim$1980 and 1990--$\sim$2005) are shifted, because neural
networks fell out of favour within AI at different times. If an exam asks
for dates, check which history the question means.
\end{notebox}

\section{Chapter Checklist}

What you should take away from this chapter:
\begin{itemize}
  \item what machine learning and deep learning are, and that learning
    \emph{representations} (increasingly abstract features) is what makes
    deep learning different;
  \item the main milestones and their order: McCulloch--Pitts (1943), Hebb
    (1949), perceptron (1957), Minsky \& Papert and XOR (1969), backprop
    (Werbos 1974; Rumelhart et al.\ 1986), universal approximation (1989),
    Hopfield networks (1982), vanishing gradients (1991), LSTM (1997), CNNs
    (Fukushima 1979; LeCun 1989), ImageNet (2009), the deep learning boom
    around 2010--2012;
  \item the two AI winters for neural networks and why they ended;
  \item the four reasons for the success since 2012: fast computers, GPUs,
    large data sets, algorithmic improvements;
  \item the four main network types and the kind of input each is made for.
\end{itemize}

% ==========================================================
%  CHAPTER 2 — BIOLOGICAL NEURAL NETWORKS   (outline)
%  Slides for this chapter are not in the course folder yet.
% ==========================================================
\chapter{Biological Neural Networks}
\label{ch:bio}

Artificial neural networks borrowed their name and their basic building
block from the brain. This chapter looks at the original: how a biological
neuron works, how neurons communicate, and which of these features survive
in the artificial neuron of \cref{ch:basic} --- and which do not.

\planned{The slides of this chapter are not yet in the course folder. The
section list below is the expected structure; adapt it to the slides once
they are available.}

\section{The Biological Neuron}
\planned{Cell body, dendrites, axon and synapses; resting potential and
action potential (``spike''); excitatory and inhibitory synapses; synaptic
strength as the biological counterpart of a weight.}

\section{Learning in Biological Networks}
\planned{Synaptic plasticity; Hebb's rule (``cells that fire together wire
together''); connection to Hebbian learning in Amari--Hopfield networks
(\cref{sec:hopfield}).}

\section{From Biological to Artificial Neurons}
\planned{The abstraction: weighted sum of inputs, activation function,
output; McCulloch--Pitts neuron; what is lost (spike timing, dendritic
computation, energy efficiency) and why the analogy should not be taken too
literally.}

\section{Chapter Checklist}
\planned{To be written together with the chapter.}


% ==========================================================
%  CHAPTER 3 — NOTATION
%  The lecture's own notation slides are not in the folder yet; this
%  chapter collects the notation as it is used on the Chapter 4 slides.
% ==========================================================
\chapter{Notation}
\label{ch:notation}

Deep learning papers, textbooks and lectures all use slightly different
symbols, and a lot of confusion in exams comes from mixing them up. This
chapter collects the notation \emph{as it is used on the Chapter~4 slides} of
this lecture, so that the rest of these notes reads consistently. The
lecture has its own notation chapter; once its slides are available, this
chapter should be checked against them.

\section{Data}
\label{sec:notation-data}

In supervised learning we are given $N$ examples (\emph{samples}, data
points), each consisting of an input vector with $D$ \emph{features} and a
label. The sample index is written as a \textbf{superscript}, so that
subscripts stay free for feature (or unit) indices:
\[
  \bx^n \in \R^D \quad\text{(input of sample } n\text{)}, \qquad
  y^n \quad\text{(its label)}, \qquad n = 1, \dots, N.
\]
The inputs are collected in the \emph{data matrix}
$\bX = (\bx^1, \dots, \bx^N)$, where the lecture's convention is that
\textbf{samples are rows} and \textbf{features are columns}, so
$\bX \in \R^{N \times D}$ and $x_{nd}$ is feature $d$ of sample $n$. The
slides' example: in a data table whose second row is $\bx^2 = (0.15, 0.14)$,
the single entry $x_{21} = 0.15$ is the first feature of the second sample;
$\bx_{\cdot 1}$ is the column of all values of feature~1. Scalar labels are
collected in $\by = (y^1, \dots, y^N)\T$.

For classification with $K$ classes, a label is written as a
\emph{one-hot vector} $\by \in \{0,1\}^K$ (\cref{sec:softmax}), and the label
vectors of all samples form the label matrix $\bY$; $y^n_k$ (or $y_{nk}$) is
entry $k$ of the label of sample $n$.

\section{Models, Losses and Learning}

\begin{itemize}
  \item $g(\bx; \bw)$ --- the network (model) as a function of the input
    $\bx$, with all its parameters collected in $\bw$ (or a weight matrix
    $\bW$, or several weight matrices $\bW\lay{1}, \dots, \bW\lay{L}$).
    A prediction is $\hat y = g(\bx; \bw)$.
  \item $\loss(y, g(\bx; \bw))$ --- the \textbf{loss function} for a single
    sample: how bad is the prediction $g(\bx; \bw)$ if the true label is $y$?
  \item $\Remp(\bw, \bX, \by)$ --- the \textbf{empirical error} (empirical
    risk): the loss summed or averaged over the training set,
    $\Remp = \frac{1}{N}\sum_n \loss(y^n, g(\bx^n; \bw))$. Training means
    minimizing $\Remp$ with respect to $\bw$.
  \item $\eta$ --- the \textbf{learning rate} (step size) of gradient
    descent: $\bw^{\text{new}} = \bw^{\text{old}} - \eta\, \grad_{\bw}
    \Remp$.
\end{itemize}

\begin{notebox}[title={Loss $\loss$ versus number of layers $L$}]
The slides use the same letter $L$ for the loss function $L(y, g(\bx;\bw))$
and for the number of layers $L$ (and $L_\nu$ for the set of units in layer
$\nu$). These notes write the loss as $\loss$ (calligraphic) to keep the two
apart. When reading the slides, the context always tells which one is meant.
\end{notebox}

\section{Units and Weights}
\label{sec:notation-units}

A neural network consists of \emph{units} (artificial neurons) connected by
weighted, directed connections.

\begin{defbox}[title={Unit-wise notation (single neurons, MLPs)}]
\begin{itemize}
  \item $w_{ij}$ --- the \textbf{weight} (connection strength, synaptic
    weight) \textbf{from unit $j$ to unit $i$}. Note the order: the
    \emph{target} comes first.
  \item $s_i = \sum_j w_{ij} a_j + b_i$ --- the \textbf{network input} or
    \textbf{pre-activation} of unit $i$; $b_i$ is its bias. Alternatively the
    bias is treated as a weight $w_{i0}$ from a constant unit $a_0 = 1$, and
    the sum runs from $j = 0$.
  \item $f$ --- the \textbf{activation function} (also \emph{transfer
    function}); it may differ between units.
  \item $a_i = f(s_i)$ --- the \textbf{activation} (activity, output) of unit
    $i$. For input units, $a_i = x_i$.
  \item Units are numbered consecutively through the whole network: the $D$
    input units are $1, \dots, D$, and if the network has $Q$ units and $K$
    outputs, the output units are $Q - K + 1, \dots, Q$.
  \item $\delta_i = \pd{\loss}{s_i}$ --- the \textbf{delta error} of unit $i$
    (\cref{sec:backprop}).
\end{itemize}
\end{defbox}

\section{Layer-Wise Notation (Deep Networks)}
\label{sec:notation-layers}

For deep networks it is easier to number the \emph{layers} and to write
everything as vectors and matrices. Layers are indexed by a superscript in
square brackets.

\begin{defbox}[title={Layer-wise notation (deep feed-forward networks)}]
\begin{itemize}
  \item $\ba\lay{0} = \bx$ --- the input is the ``activation of layer 0''.
  \item $\bW\lay{l}$ --- the weight matrix connecting layer $l-1$ to layer
    $l$; entry $w\lay{l}_{ij}$ goes from unit $j$ in layer $l-1$ to unit $i$
    in layer $l$. If layer $l$ has $n_l$ units, $\bW\lay{l} \in
    \R^{n_l \times n_{l-1}}$. Bias vector: $\bb\lay{l}$.
  \item $\bs\lay{l} = \bW\lay{l} \ba\lay{l-1} + \bb\lay{l}$ --- the
    pre-activations of layer $l$.
  \item $\ba\lay{l} = f(\bs\lay{l})$ --- the activations of layer $l$, with
    $f$ applied element-wise.
  \item $\byhat = \ba\lay{L}$ --- the output of the last layer $L$ (the
    output activation is often written $\sigma$ to stress that it may be
    different from the hidden-layer $f$).
  \item $\bdelta\lay{l} = \bigl(\pd{\loss}{\bs\lay{l}}\bigr)\T$ --- the delta
    errors of layer $l$ (\cref{sec:backprop-dnn}).
\end{itemize}
\end{defbox}

\begin{table}[htbp]
\centering
\small
\begin{tabular}{@{}lll@{}}
\toprule
\textbf{Quantity} & \textbf{Unit-wise (MLP)} & \textbf{Layer-wise (DNN)} \\
\midrule
input              & $a_i = x_i$, $i \le D$   & $\ba\lay{0} = \bx$ \\
weight             & $w_{ij}$ (from $j$ to $i$) & $\bW\lay{l}$, entries $w\lay{l}_{ij}$ \\
pre-activation     & $s_i = \sum_j w_{ij} a_j + b_i$ & $\bs\lay{l} = \bW\lay{l}\ba\lay{l-1} + \bb\lay{l}$ \\
activation         & $a_i = f(s_i)$          & $\ba\lay{l} = f(\bs\lay{l})$ \\
output             & $a_k$, $Q-K+1 \le k \le Q$ & $\byhat = \ba\lay{L}$ \\
delta error        & $\delta_i = \partial\loss/\partial s_i$ & $\bdelta\lay{l} = (\partial\loss/\partial\bs\lay{l})\T$ \\
\bottomrule
\end{tabular}
\caption{The two notations side by side. The unit-wise notation is used for
single-layer networks and MLPs (Sections~4.1--4.6), the layer-wise notation
for deep feed-forward networks (Section~4.7).}
\label{tab:notation}
\end{table}

\section{Further Conventions}

\begin{itemize}
  \item Vectors are column vectors and written in bold italic ($\bx$, $\bw$);
    matrices are bold capitals ($\bW$, $\bX$). Transposition: $\bw\T$.
  \item $\sigma(x) = 1/(1 + e^{-x})$ is the logistic sigmoid; $\softmax$
    is defined in \cref{sec:softmax}.
  \item $\grad_{\bw} \Remp$ is the gradient (a vector of the same shape as
    $\bw$); $\hess_{\bw} \Remp$ the Hessian (\cref{sec:hessian}).
  \item $\diag(\bm{v})$ is the diagonal matrix with $\bm{v}$ on the diagonal;
    $\odot$ is the element-wise product.
\end{itemize}

% ==========================================================
%  CHAPTER 4 — BASIC NEURAL NETWORK ARCHITECTURES
%  Sources: "TEX Chapter4 - Basics" (4.1--4.4, Hopfield) and
%           "Chapter4 - FNN" (4.5--4.10)
% ==========================================================
\chapter{Basic Neural Network Architectures}
\label{ch:basic}

This is the central chapter of the course. It starts with the simplest
possible network --- a single neuron that computes a weighted sum --- and
adds one idea at a time: a threshold (perceptron), a sigmoid (logistic
regression), a softmax (multi-class classification), a hidden layer (MLP),
many hidden layers (deep networks). Along the way it develops the two tools
that every later chapter relies on: \emph{maximum likelihood} to decide
which loss to minimize, and \emph{gradient descent with backpropagation} to
minimize it.

The section numbers follow the lecture: Sections 4.1--4.4 cover
\emph{single-layer} networks, 4.5 shows why they are not enough, 4.6--4.7
cover multi-layer and deep networks with backpropagation, 4.8--4.9 the
Jacobian and the Hessian, and 4.10 the practical tricks of the trade. An
excursion to Amari--Hopfield networks, which closed the single-layer lecture,
comes at the end (\cref{sec:hopfield}).

\paragraph{Units, weights and layers}
An artificial neural network is a set of \emph{units} connected by directed,
weighted connections. The weight $w_{ij}$ is the connection strength
(``synaptic weight'') from unit $j$ to unit $i$. Units are organized in
\emph{layers}: an \emph{input} layer that receives the data, one or more
\emph{hidden} layers, and an \emph{output} layer (\cref{fig:three-layer}).
Learning means adapting the weights so that the network computes something
useful. The notation is collected in \cref{ch:notation}.

\begin{figure}[htbp]
\centering
\begin{tikzpicture}[node distance=12mm]
  % input
  \node[inunit] (x1) at (-0.7,0) {};
  \node[inunit] (x2) at (0.7,0) {};
  % hidden
  \foreach \i/\xx in {1/-2.1,2/-0.7,3/0.7,4/2.1}
    \node[unit] (h\i) at (\xx,1.6) {};
  % output
  \node[outunit] (o) at (0,3.2) {};
  \foreach \i in {1,...,4} {
    \draw[conn] (x1) -- (h\i);
    \draw[conn] (x2) -- (h\i);
    \draw[conn] (h\i) -- (o);
  }
  \node[font=\small, anchor=east] at (-3.0,0) {input};
  \node[font=\small, anchor=east] at (-3.0,1.6) {hidden};
  \node[font=\small, anchor=east] at (-3.0,3.2) {output};
  % zoom on a single weight
  \begin{scope}[xshift=6.3cm]
    \node[inunit] (j) at (0.6,0.4) {$j$};
    \node[unit] (i) at (0.6,2.8) {$i$};
    \draw[conn, thick] (j) -- node[right, font=\small] {$w_{ij}$} (i);
    \node[font=\scriptsize, align=left, anchor=west] at (1.4,1.6)
      {weight \emph{from} unit $j$\\ \emph{to} unit $i$};
  \end{scope}
\end{tikzpicture}
\caption{Left: a three-layer network with an input, a hidden and an output
layer. Right: the weight $w_{ij}$ connects unit $j$ to unit $i$ --- the
\emph{target} index comes first.}
\label{fig:three-layer}
\end{figure}

\paragraph{Data for supervised learning}
Throughout the chapter we have a supervised data set of $N$ samples. In
\emph{regression} the labels are real numbers, $y^n \in \R$; in
\emph{classification} they are class labels. The inputs form the data
matrix $\bX \in \R^{N \times D}$ (one sample $\bx^n$ per row), and the labels
form the vector $\by = (y^1, \dots, y^N)\T$.


% ----------------------------------------------------------
\section{Linear Model: The Linear Neuron}
\label{sec:linear-neuron}
% ----------------------------------------------------------

The simplest network has no hidden layer at all: the inputs are directly
connected to a single output unit, which outputs the weighted sum of its
inputs.

\begin{defbox}[title={Linear neuron}]
For an input $\bx \in \R^D$ and adaptive weights $\bw \in \R^D$, the linear
neuron computes the activation
\[
  g(\bx; \bw) = a = \bw\T\bx = \sum_{d=1}^{D} w_d x_d \;(+\, w_0).
\]
A bias $w_0$ can be included by appending a constant feature $x_0 = 1$ to
every input, so we can always write $\bw\T\bx$. The linear neuron is exactly
the model of \textbf{linear regression}.
\end{defbox}

\subsection{From Gaussian Noise to the Squared Loss}

Which weights are ``good''? The lecture answers this with a probabilistic
model, and it is worth understanding the argument, because the same recipe
gives us the loss functions for every other network in this chapter.

Assume the label is a deterministic function of the input plus Gaussian
noise:
\[
  y = g(\bx; \bw) + \epsilon, \qquad \epsilon \sim \Normal(0, \sigma^2).
\]
Then the label itself is Gaussian around the network output,
$p(y \given \bx, \bw, \sigma) = \Normal(y \given g(\bx;\bw), \sigma^2)$,
and the best prediction (the conditional mean) is the network output:
$\E[y \given \bx] = \int y\, p(y \given \bx) \dd y = g(\bx; \bw)$.

If the $N$ data points are drawn independently, their joint probability ---
the \emph{likelihood} of the weights --- is the product
\[
  p(\by \given \bX, \bw, \sigma)
  = \prod_{n=1}^{N} \Normal(y^n \given \bw\T\bx^n, \sigma^2)
  = \prod_{n=1}^{N} \frac{1}{\sqrt{2\pi\sigma^2}}
    \exp\!\Bigl(-\frac{(y^n - \bw\T\bx^n)^2}{2\sigma^2}\Bigr).
\]
Taking the logarithm (which does not move the location of the maximum)
turns the product into a sum:
\[
  \log p(\by \given \bX, \bw, \sigma)
  = -\frac{N}{2}\log(2\pi\sigma^2) - \frac{1}{\sigma^2}\, S(\bw),
  \qquad
  S(\bw) = \frac12 \sum_{n=1}^{N} \bigl(\underbrace{y^n - \bw\T\bx^n}_{=:\,\epsilon^n}\bigr)^2 .
\]
The first term does not depend on $\bw$, and $\sigma^2$ only scales the
second. So maximizing the likelihood is the same as minimizing $S(\bw)$, the
(halved) \emph{sum of squared residuals} $\epsilon^n$.

\begin{thmbox}[title={Maximum likelihood $\Leftrightarrow$ least squares}]
Under the assumption of additive Gaussian noise, maximizing the likelihood
of the data is equivalent to minimizing the \textbf{squared loss}
\begin{align*}
  \loss(y, g(\bx;\bw)) &= \tfrac12\bigl(y - g(\bx;\bw)\bigr)^2, \\
  \Remp(\bw) = S(\bw) &= \tfrac12 \sum_{n=1}^{N} (y^n - \bw\T\bx^n)^2
  = \tfrac12 (\by - \bX\bw)\T(\by - \bX\bw).
\end{align*}
A candidate weight vector $\tilde\bw$ with a low sum of squared residuals
$S(\tilde\bw)$ has a high likelihood.
\end{thmbox}

\begin{notebox}[title={$\Normal(0,\sigma)$ on the slides}]
The slides write the noise as $\epsilon \sim \Normal(0, \sigma)$, with
$\sigma$ the \emph{standard deviation}. These notes write the second
argument as the variance, $\Normal(0,\sigma^2)$. The formulas are the same
either way --- just do not plug $\sigma$ in where $\sigma^2$ belongs.
\end{notebox}

\subsection{Learning: Empirical Risk Minimization}

Learning now means finding the weights that minimize the error on the
training set,
\[
  \what = \argmin_{\tilde\bw} S(\tilde\bw),
\]
an approach called \textbf{empirical risk minimization} (ERM) with
$\Remp(\tilde\bw, \bX, \by) = S(\tilde\bw)$. Equivalently, we minimize the
\emph{negative log-likelihood}
$-\log p(\by \given \bX, \bw, \sigma) = \frac{N}{2}\log(2\pi\sigma^2) +
\frac{1}{\sigma^2} S(\bw)$. Its gradient is
\[
  \grad_{\bw}\bigl(-\log p(\by \given \bX, \bw, \sigma)\bigr)
  = \frac{1}{\sigma^2} \sum_{n=1}^{N} (\bw\T\bx^n - y^n)\,\bx^n
  = \frac{1}{\sigma^2}\bigl(\bX\T\bX\bw - \bX\T\by\bigr).
\]
Note the structure of each summand: \emph{(prediction $-$ target) $\times$
input}. We will meet it again for every single-layer network.

There are two ways to find the minimum: solve for it directly, or walk
towards it step by step.

\subsection{Closed-Form Solution: The Normal Equations}

Setting the gradient to zero gives the \emph{normal equations}:
\[
  0 = \bX\T\bX\bw - \bX\T\by
  \quad\Longleftrightarrow\quad
  \bX\T\bX\,\bw = \bX\T\by
  \quad\Longrightarrow\quad
  \what = (\bX\T\bX)\inv\bX\T\by .
\]

\begin{thmbox}[title={Least-squares solution}]
If $\bX\T\bX$ has full rank (is invertible), the unique minimizer of the
squared error --- and the maximum-likelihood estimator under Gaussian noise
--- is
\[
  \what = (\bX\T\bX)\inv\bX\T\by = \bX^{+}\by,
  \qquad \bX^{+} = (\bX\T\bX)\inv\bX\T,
\]
where $\bX^{+}$ is the \textbf{Moore--Penrose pseudo-inverse} of $\bX$.
For centered data, $\frac1N\bX\T\bX$ is the estimate of the covariance
matrix of the features.
\end{thmbox}

\begin{exbox}[title={Fitting a line through three points}]
Data: $x = 0, 1, 2$ with labels $\by = (1, 3, 4)\T$. To fit an intercept, we
append the constant feature, so each row of $\bX$ is $(1, x^n)$:
\[
  \bX = \begin{pmatrix} 1 & 0 \\ 1 & 1 \\ 1 & 2 \end{pmatrix}, \quad
  \bX\T\bX = \begin{pmatrix} 3 & 3 \\ 3 & 5 \end{pmatrix}, \quad
  \bX\T\by = \begin{pmatrix} 8 \\ 11 \end{pmatrix}, \quad
  (\bX\T\bX)\inv = \frac16\begin{pmatrix} 5 & -3 \\ -3 & 3 \end{pmatrix}.
\]
Hence $\what = \frac16(40 - 33,\; -24 + 33)\T = (7/6,\; 3/2)\T \approx
(1.167, 1.5)\T$, i.e.\ the line $\hat y = 1.167 + 1.5\,x$. The predictions are
$(1.167, 2.667, 4.167)$, the residuals $(-0.167, 0.333, -0.167)$, and
$S(\what) = \frac12(0.0278 + 0.1111 + 0.0278) = 0.083$. No other line through
these points has a smaller squared error.
\end{exbox}

\subsection{Gradient Descent and Stochastic Gradient Descent}

Instead of solving the normal equations, we can start from an initial
guess and improve it repeatedly by stepping against the gradient:
\[
  \bw^{\text{new}} = \bw^{\text{old}} - \eta\, \grad_{\bw} S(\bw)\big|_{\bw^{\text{old}}},
  \qquad
  \grad_{\bw} S(\bw) = \sum_{n=1}^{N} (\bw\T\bx^n - y^n)\,\bx^n .
\]
This is \textbf{gradient descent}. The \emph{learning rate} $\eta$ is
critical: too small and the algorithm crawls, too large and it overshoots
and diverges. For least squares, gradient descent works well because the
problem is \emph{convex}: there are no bad local minima to get stuck in.

\begin{exbox}[title={One gradient step on the three points}]
Start at $\bw = (0,0)\T$, where $S = \frac12(1 + 9 + 16) = 13$. The gradient
is $\bX\T(\bX\bw - \by) = -\bX\T\by = (-8, -11)\T$. With $\eta = 0.1$ the new
weights are $\bw = (0.8, 1.1)\T$, the predictions $(0.8, 1.9, 3.0)$, and the
error drops to $S = \frac12(0.04 + 1.21 + 1.00) = 1.125$. Repeating the step
converges to the closed-form solution $(1.167, 1.5)\T$ with $S = 0.083$.
\end{exbox}

For large data sets, computing the full gradient over all $N$ samples for
every single step is expensive. Here a simple observation helps: the scaled
gradient is a \emph{mean} over data points,
\[
  \frac1N \grad_{\bw} S(\bw) = \frac1N\sum_{n=1}^{N} (\bw\T\bx^n - y^n)\,\bx^n,
\]
and a mean can be estimated from a random subset. Drawing a
\textbf{mini-batch}
\[
  B = \{\bx^{n_1}, \dots, \bx^{n_{|B|}}\} \subset \{\bx^1, \dots, \bx^N\}
\]
gives
\[
  \frac1N \grad_{\bw} S(\bw) \approx
  \frac{1}{|B|} \sum_{b=1}^{|B|} (\bw\T\bx^{n_b} - y^{n_b})\,\bx^{n_b},
\]
and using this noisy but cheap estimate in the update is \textbf{stochastic
gradient descent} (SGD). It is the workhorse of all deep learning
(\cref{sec:tricks}).

\subsection{The Trained Network}

The goal of learning is a trained network that makes predictions for new
inputs, $g(\bx^{\text{new}}; \bw) = \bw\T\bx^{\text{new}}$. In the slides'
example, the weights $\bw = (0.1, -1.0, 0.4)\T$ and the new input
$\bx^{\text{new}} = (1.0, -0.3, 0)\T$ give the prediction
\[
  a = 0.1\cdot 1.0 + (-1.0)\cdot(-0.3) + 0.4 \cdot 0 = 0.4 .
\]

\begin{intbox}
The linear neuron already contains the whole pipeline of this course in
miniature: an \emph{architecture} (weighted sum), an \emph{objective}
(empirical error, motivated by maximum likelihood), \emph{learning}
(minimizing the objective --- in closed form or by gradient descent), and
the \emph{goal}: a trained network that predicts well on new data. Every
network that follows changes one of these pieces.
\end{intbox}


% ----------------------------------------------------------
\section{Perceptron}
\label{sec:perceptron}
% ----------------------------------------------------------

Now the labels are not real numbers but two classes, coded as
$y \in \{-1, +1\}$: diseased vs.\ healthy, spam vs.\ normal email, win or
loss. We could simply run linear regression on the labels $\pm 1$, but this is
a poor fit to the problem: the squared loss also punishes points that are
classified correctly with a large margin (a prediction of $+5$ for a label
$+1$ is ``wrong'' by 4), and the output is a real number, not a class.

\begin{defbox}[title={Perceptron}]
The perceptron is a linear neuron followed by the sign function:
\[
  g(\bx; \bw) = a = \sign(\bw\T\bx), \qquad
  \sign(z) = \begin{cases} +1 & z > 0, \\ \phantom{+}0 & z = 0, \\ -1 & z < 0. \end{cases}
\]
It predicts class $+1$ on one side of the hyperplane $\bw\T\bx = 0$ and class
$-1$ on the other. The weight vector $\bw$ is the \emph{normal vector} of this
decision boundary and points to the $+1$ side.
\end{defbox}

\subsection{Perceptron Loss and Learning Rule}

A sample is classified correctly if $y$ and $\bw\T\bx$ have the same sign,
i.e.\ if $y\,\bw\T\bx > 0$. The \textbf{perceptron loss} is zero for correct
samples and grows linearly with how far a misclassified sample is on the
wrong side:
\[
  \loss(y, \bw\T\bx) = \max\bigl(0,\; -y\,\bw\T\bx\bigr),
  \qquad
  \Remp(\by, \bX, \bw) = -\sum_{n \in \mathcal{M}} y^n\, \bw\T\bx^n,
\]
where $\mathcal{M}$ is the set of \textbf{misclassified} samples (the slides
write the sum as running over all $n$ with $a^n y^n = -1$). Each term is
positive, so $\Remp \ge 0$, with equality exactly when all samples are
classified correctly.

The gradient is $\grad_{\bw}\Remp = -\sum_{n\in\mathcal{M}} y^n\bx^n$, and a
gradient-descent step gives the \textbf{perceptron learning rule}:
\[
  \Delta\bw = \eta \sum_{n \in \mathcal{M}} y^n\, \bx^n
  \quad\text{(whole data set)},
  \qquad
  \Delta\bw = \eta\, y\, \bx \;\text{ if } \bx \text{ is misclassified}
  \quad\text{(single sample)}.
\]
In the online version, we go through the samples one by one and change the
weights only when the current sample is misclassified.

\begin{notebox}[title={Sign conventions on the slides}]
The slides write the perceptron loss as ``$y(\bw\T\bx)$ for misclassified
samples''; with that sign it would be \emph{negative}. The loss must be
$-y\,\bw\T\bx$ (as in the slides' empirical error, which has the minus sign).
Also, with $\sign(0) = 0$, a sample exactly on the boundary would count as
correct, and a perceptron started at $\bw = \bm{0}$ would never move. In
practice --- and in the convergence proof --- a sample counts as
misclassified if $y\,\bw\T\bx \le 0$.
\end{notebox}

\subsection{Geometric Interpretation}

The learning rule has a nice geometric meaning (\cref{fig:perceptron-rule}).
If a sample $\bx$ of class $+1$ lies on the wrong side, then $\bw\T\bx < 0$:
the angle between $\bw$ and $\bx$ is larger than $90^\circ$. Adding $\eta\bx$
rotates $\bw$ \emph{towards} $\bx$, and the boundary rotates with it. For a
misclassified sample of class $-1$, the update $\bw - \eta\bx$ rotates $\bw$
\emph{away} from $\bx$. Either way, $y\,\bw\T\bx$ increases by
$\eta\norm{\bx}^2$: the update makes this particular sample ``more
correct''. It may, however, make other samples wrong.

\begin{figure}[htbp]
\centering
\begin{tikzpicture}[scale=1.25,
    cpos/.style={font=\small\bfseries, figgreen},
    cneg/.style={font=\small\bfseries, figred}]
  \draw[gray!60, -{Stealth[length=1.6mm]}] (-2.2,0) -- (2.3,0) node[right, font=\scriptsize] {$x_1$};
  \draw[gray!60, -{Stealth[length=1.6mm]}] (0,-2.0) -- (0,2.1) node[above, font=\scriptsize] {$x_2$};
  % old and new boundary
  \draw[accent, thick] (-0.632,1.897) -- (0.632,-1.897)
    node[below, font=\scriptsize, accent] {old boundary};
  \draw[accent, thick, dashed] (1.728,-1.008) -- (-1.728,1.008)
    node[left, font=\scriptsize, accent, align=right] {new\\boundary};
  % points
  \node[cpos] at (1.1,1.75) {$+$};
  \node[cpos] at (1.9,0.6) {$+$};
  \node[cneg] at (-1.4,-0.6) {$-$};
  \node[cneg] at (-0.4,-1.4) {$-$};
  \node[cneg] at (0.2,-1.6) {$-$};
  \node[cpos, draw=figgreen, circle, inner sep=0.5pt] (xx) at (-1.0,1.6) {$+$};
  \node[font=\scriptsize, figgreen!70!black, anchor=east] at (-1.15,1.6) {$\bx$ (misclassified)};
  % vectors
  \draw[-{Stealth[length=2.2mm]}, very thick, accent] (0,0) -- (1.2,0.4)
    node[below right, font=\small] {$\bw^{\text{old}}$};
  \draw[-{Stealth[length=2mm]}, thick, figgreen, dashed] (1.2,0.4) -- (0.7,1.2)
    node[midway, right, font=\small] {$\eta\bx$};
  \draw[-{Stealth[length=2.2mm]}, very thick, figblue] (0,0) -- (0.7,1.2)
    node[above, font=\small] {$\bw^{\text{new}}$};
\end{tikzpicture}
\caption{One perceptron update. The $+$ sample $\bx$ lies on the negative side
of the old boundary (solid). Adding $\eta\bx$ to the weight vector rotates
$\bw$ towards $\bx$; the new boundary (dashed) is perpendicular to
$\bw^{\text{new}}$ and now classifies all points correctly. Here
$\bw^{\text{old}} = (1.2, 0.4)$, $\bx = (-1, 1.6)$, $\eta = 0.5$.}
\label{fig:perceptron-rule}
\end{figure}

\begin{exbox}[title={Running the perceptron by hand}]
Four samples in 2D, each extended by a constant feature 1 for the bias:
$\bx^1 = (2, 1, 1)$, $\bx^2 = (0, 2, 1)$ with $y = +1$ and
$\bx^3 = (-1, 0, 1)$, $\bx^4 = (1, -2, 1)$ with $y = -1$. Start with
$\bw = (0, 0, 0)$, $\eta = 1$, and update whenever $y\,\bw\T\bx \le 0$.

\medskip
\centering
\begin{tabular}{@{}cccrl@{}}
\toprule
epoch & sample & $\bw$ before & $y\,\bw\T\bx$ & action \\
\midrule
1 & $\bx^1$ & $(0, 0, 0)$ & $0$  & update: $\bw = (2, 1, 1)$ \\
1 & $\bx^2$ & $(2, 1, 1)$ & $3$  & correct \\
1 & $\bx^3$ & $(2, 1, 1)$ & $1$  & correct \\
1 & $\bx^4$ & $(2, 1, 1)$ & $-1$ & update: $\bw = (2,1,1) - (1,-2,1) = (1, 3, 0)$ \\
2 & all     & $(1, 3, 0)$ & $5, 6, 1, 5$ & all correct --- stop \\
\bottomrule
\end{tabular}

\raggedright
\smallskip
After two updates the perceptron has found the separating line
$x_1 + 3x_2 = 0$. A different order of the samples would have found a
different line.
\end{exbox}

\subsection{Convergence for Linearly Separable Data}

Does the learning rule always find a solution? If one exists, yes.

\begin{thmbox}[title={Perceptron convergence theorem}, unbreakable]
If the training data are \textbf{linearly separable} --- there is a $\what$
with $y^n\,\what\T\bx^n > 0$ for all $n$ --- then the perceptron learning rule
finds a separating weight vector after a \emph{finite} number of updates.
With margin $\gamma = \min_n y^n\,\what\T\bx^n$ and $R = \max_n\norm{\bx^n}$,
the number of updates is at most $\norm{\what}^2 R^2 / \gamma^2$.
\end{thmbox}

\paragraph{Proof idea} Start with $\bw = \bm{0}$ and update with $\eta = 1$
(the step size does not matter when starting from zero). If sample $n$ has
been misclassified $\tau_n$ times, then after $\tau = \sum_n \tau_n$ updates
$\bw = \sum_n \tau_n\, y^n\bx^n$. Now compare two quantities:
\begin{enumerate}
  \item \emph{The projection onto $\what$ grows linearly.} Every update adds
    at least $\gamma$, so
    \[
      \what\T\bw = \sum_n \tau_n\, \what\T\bx^n y^n \ge \tau\gamma .
    \]
  \item \emph{The length grows only like $\sqrt\tau$.} An update on a
    misclassified sample changes the squared length by
    $\norm{\bw + y^n\bx^n}^2 - \norm{\bw}^2 = \norm{\bx^n}^2 + 2y^n\bw\T\bx^n
    \le R^2$, because the last term is $\le 0$ for a misclassified sample.
    So $\norm{\bw}^2 \le \tau R^2$.
\end{enumerate}
By Cauchy--Schwarz, $\tau\gamma \le \what\T\bw \le \norm{\what}\norm{\bw}
\le \norm{\what}\sqrt{\tau}R$. The left side grows like $\tau$, the right
side only like $\sqrt\tau$ --- for large $\tau$ the two would be
incompatible. Hence $\tau \le \norm{\what}^2R^2/\gamma^2$. \hfill$\square$

\begin{notebox}[title={What the theorem does \emph{not} say}]
\begin{itemize}
  \item \textbf{The solution is not unique.} Any separating hyperplane has
    zero perceptron loss. Which one the algorithm finds depends on the
    initialization and on the order of the samples.
  \item \textbf{Non-separable data:} if no separating hyperplane exists, the
    perceptron never stops updating --- the weights keep cycling, and the
    algorithm does not converge.
  \item The bound can be large: a small margin $\gamma$ means many updates.
\end{itemize}
\end{notebox}


% ----------------------------------------------------------
\section{Linear Neuron with Sigmoid Activation: Logistic Regression}
\label{sec:logreg}
% ----------------------------------------------------------

Again we have two classes, now coded as $y \in \{0, 1\}$. Instead of a hard
decision we would like the network to output a \emph{probability} for
class~1. The trick: squash the weighted sum into $(0,1)$.

\begin{defbox}[title={Logistic neuron (logistic regression)}]
\[
  g(\bx; \bw) = a = \sigma(\bw\T\bx) = \frac{1}{1 + e^{-\bw\T\bx}},
\]
where $\sigma$ is the \textbf{logistic sigmoid}. The output is interpreted as
$p(y = 1 \given \bx; \bw)$. This is the model of logistic regression.
\end{defbox}

The sigmoid (\cref{fig:sigmoid}) is $0.5$ at zero, approaches $1$ for large
positive and $0$ for large negative inputs, and is symmetric:
$\sigma(-z) = 1 - \sigma(z)$. Its derivative has a particularly convenient
form, which we will use over and over:
\[
  \sigma'(z) = \frac{\dd}{\dd z}\frac{1}{1 + e^{-z}}
  = \frac{e^{-z}}{(1+e^{-z})^2} = \sigma(z)\bigl(1 - \sigma(z)\bigr).
\]
It is largest at $z = 0$, where $\sigma'(0) = 0.5 \cdot 0.5 = 0.25$, and
goes to zero for large $|z|$ --- the sigmoid \emph{saturates}.

\begin{figure}[htbp]
\centering
\begin{tikzpicture}
\begin{axis}[dlbt, width=0.62\linewidth, height=4.6cm,
    xmin=-6.2, xmax=6.2, ymin=-0.05, ymax=1.12,
    xtick={-6,-4,...,6}, ytick={0.25,0.5,1},
    legend style={font=\scriptsize, draw=none, fill=white, fill opacity=0.85,
                  at={(0.02,0.97)}, anchor=north west}]
  \addplot[very thick, accent, domain=-6:6] {1/(1+exp(-x))};
  \addlegendentry{$\sigma(z)$}
  \addplot[very thick, figamber, dashed, domain=-6:6] {exp(-x)/(1+exp(-x))^2};
  \addlegendentry{$\sigma'(z) = \sigma(z)(1-\sigma(z))$}
\end{axis}
\end{tikzpicture}
\caption{The logistic sigmoid and its derivative. The derivative never
exceeds $0.25$ --- a fact that becomes important for the vanishing gradient
problem (\cref{sec:vgp}).}
\label{fig:sigmoid}
\end{figure}

\subsection{Likelihood and Binary Cross-Entropy}

The recipe is the same as for the linear neuron: write down the probability
of the data and maximize it. The network says
$p(y = 1 \given \bx; \bw) = g(\bx;\bw)$ and $p(y = 0 \given \bx; \bw) = 1 -
g(\bx;\bw)$, which can be written compactly as a Bernoulli distribution
\[
  p(y \given \bx; \bw) = g(\bx;\bw)^{y}\,\bigl(1 - g(\bx;\bw)\bigr)^{1-y}.
\]
(For $y = 1$ the second factor is 1, for $y = 0$ the first.) For independent
samples, the likelihood and the negative log-likelihood are
\[
  p(\by \given \bX, \bw) = \prod_{n=1}^{N} g(\bx^n;\bw)^{y^n}\bigl(1 - g(\bx^n;\bw)\bigr)^{1-y^n},
\]
\[
  -\log p(\by \given \bX, \bw)
  = -\sum_{n=1}^{N} \Bigl( y^n \log g(\bx^n;\bw) + (1 - y^n)\log\bigl(1 - g(\bx^n;\bw)\bigr) \Bigr).
\]

\begin{defbox}[title={Binary cross-entropy}]
The loss
\[
  \loss(y, g(\bx;\bw)) = -\Bigl(y \log g(\bx;\bw) + (1-y)\log\bigl(1 - g(\bx;\bw)\bigr)\Bigr)
\]
is called \textbf{binary cross-entropy}. It is the negative log-probability
that the model assigns to the correct class: small if the model is
confidently right, very large if it is confidently wrong.
\end{defbox}

\subsection{Gradient and Hessian}

Using $\sigma' = \sigma(1-\sigma)$ and writing $a^n = \sigma(\bw\T\bx^n)$:
\begin{align*}
  \grad_{\bw}\Remp
  &= -\sum_{n=1}^{N} \Bigl( \frac{y^n}{a^n} - \frac{1 - y^n}{1 - a^n} \Bigr)
     \grad_{\bw}\sigma(\bw\T\bx^n)
   = -\sum_{n=1}^{N} \Bigl( \frac{y^n}{a^n} - \frac{1 - y^n}{1 - a^n} \Bigr)
     a^n(1 - a^n)\,\bx^n \\
  &= -\sum_{n=1}^{N} \bigl( y^n(1 - a^n) - (1 - y^n)a^n \bigr)\bx^n
   = \sum_{n=1}^{N} \bigl(\sigma(\bw\T\bx^n) - y^n\bigr)\,\bx^n .
\end{align*}

\begin{thmbox}[title={Logistic regression: gradient and Hessian}]
For a single sample and for the whole data set:
\[
  \grad_{\bw}\Remp = (\sigma(\bw\T\bx) - y)\,\bx,
  \qquad
  \grad_{\bw}\Remp = \sum_{n=1}^{N}(\sigma(\bw\T\bx^n) - y^n)\,\bx^n = \bX\T(\ba - \by),
\]
\[
  \hess_{\bw}\Remp = \sigma(\bw\T\bx)\bigl(1 - \sigma(\bw\T\bx)\bigr)\,\bx\bx\T,
  \qquad
  \hess_{\bw}\Remp = \sum_{n=1}^{N} a^n(1 - a^n)\,\bx^n\bx^{n\mathsf{T}} = \bX\T\bR\bX,
\]
where $\ba = (a^1, \dots, a^N)\T$ and $\bR = \diag\bigl(a^n(1 - a^n)\bigr)$.
Since all diagonal entries of $\bR$ are positive, $\bX\T\bR\bX$ is positive
semi-definite: the binary cross-entropy is \textbf{convex} in $\bw$.
\end{thmbox}

The gradient looks exactly like the one of linear regression ---
\emph{(prediction $-$ target) $\times$ input} --- only the prediction is now
$\sigma(\bw\T\bx)$ instead of $\bw\T\bx$. This is no coincidence, and it will
save us a lot of work in backpropagation (\cref{tab:loss-activation}).

\begin{notebox}[title={The sign of the Hessian}]
One slide writes the Hessian of logistic regression as $-\bX\T\bR\bX$; that
is the Hessian of the \emph{log-likelihood}, which is maximized. For the
error $\Remp$ (negative log-likelihood), which is minimized, the Hessian is
$+\bX\T\bR\bX$, and this is the matrix that appears in the update rule
below.
\end{notebox}

\subsection{Learning}

Setting the gradient to zero gives no closed-form solution this time (the
$\sigma$ is in the way), so we have to iterate:
\begin{itemize}
  \item \textbf{gradient descent:}
    $\bw^{\text{new}} = \bw^{\text{old}} - \eta\,\grad_{\bw}\Remp(\by, \bX, \bw)|_{\bw^{\text{old}}}$;
  \item \textbf{stochastic gradient descent:} the same with a mini-batch
    estimate $\tilde\grad_{\bw}\Remp$ of the gradient;
  \item \textbf{iteratively reweighted least squares} (IRLS), the method
    typically used in statistics software: Newton's method with the Hessian,
    $\bw^{\text{new}} = \bw^{\text{old}} - \bH\inv\grad_{\bw}\Remp|_{\bw^{\text{old}}}$
    (derived below).
\end{itemize}

\begin{exbox}[title={One gradient step for logistic regression}]
Weights $\bw = (0.5, -1)\T$, one sample $\bx = (1, 2)\T$ with label $y = 1$.
\begin{itemize}
  \item Forward: $\bw\T\bx = 0.5 - 2 = -1.5$, $a = \sigma(-1.5) = 0.182$. The
    model gives the correct class only 18\,\% probability; the loss is
    $-\log 0.182 = 1.70$.
  \item Gradient: $(a - y)\bx = (0.182 - 1)(1, 2)\T = (-0.818, -1.635)\T$.
  \item Step with $\eta = 0.5$: $\bw = (0.5, -1)\T + 0.5\,(0.818, 1.635)\T
    = (0.909, -0.182)\T$.
  \item Check: $\bw\T\bx = 0.544$, $a = \sigma(0.544) = 0.633$, loss $0.458$.
\end{itemize}
One step moved the prediction from ``probably class 0'' to ``probably class
1''. Note how large the step was: the gradient is proportional to the error
$a - y$, which was big.
\end{exbox}

\subsection{Intermezzo: Optimization and Newton's Method}

Learning is an optimization problem, so it is worth placing it in the
general picture. A generic (constrained) optimization problem reads
\[
  \min_{\bx}\; f(\bx) \quad\text{s.t.}\quad g(\bx) \le 0,\;\; h(\bx) = 0,
  \qquad f: \R^n \to \R,\; g: \R^n \to \R^m,\; h: \R^n \to \R .
\]
(This is general optimization notation, not the course's ML notation.) In
neural-network training, the variable $\bx$ is our \emph{weight vector},
$f$ is the \emph{empirical error}, and there are typically \emph{no
constraints}. Optimization methods differ in what they may ask about $f$:
\begin{itemize}
  \item \textbf{black-box optimization:} only function values $f(\bx)$;
  \item \textbf{gradient-based optimization:} $f(\bx)$ and the gradient
    $\grad f(\bx)$ --- this is what deep learning uses;
  \item \textbf{second-order optimization:} additionally the Hessian
    $\hess f(\bx)$.
\end{itemize}

\paragraph{Newton's method}
To find a \emph{root} of a function, approximate it by its tangent (a
first-order Taylor expansion) and solve for the zero of the tangent:
\[
  f(x) \approx f(x_0) + (x - x_0) f'(x_0) = 0
  \quad\Longrightarrow\quad
  x = x_0 - \frac{f(x_0)}{f'(x_0)} .
\]
To find an \emph{optimum}, apply the same idea to the derivative (an
optimum is a root of $f'$). In one and in $n$ dimensions:
\[
  x^{\text{new}} = x^{\text{old}} - \frac{f'(x^{\text{old}})}{f''(x^{\text{old}})},
  \qquad
  \bx^{\text{new}} = \bx^{\text{old}} - \bigl(\hess f(\bx^{\text{old}})\bigr)\inv \grad f(\bx^{\text{old}}).
\]
Compared with gradient descent, the Hessian replaces the fixed learning rate
by a step that is adapted to the curvature in every direction.

\begin{thmbox}[title={Newton's method for linear and logistic regression}]
With $\bw^{\text{new}} = \bw^{\text{old}} - (\hess_{\bw}\Remp)\inv\grad_{\bw}\Remp$:
\begin{itemize}
  \item \textbf{Linear regression:} $\grad_{\bw}\Remp = \bX\T\bX\bw - \bX\T\by$,
    $\hess_{\bw}\Remp = \bX\T\bX$, so
    \[
      \bw^{\text{new}} = \bw^{\text{old}} - (\bX\T\bX)\inv(\bX\T\bX\bw^{\text{old}} - \bX\T\by)
      = (\bX\T\bX)\inv\bX\T\by .
    \]
    A single Newton step from \emph{any} starting point lands exactly on the
    least-squares solution --- the error is quadratic, so its second-order
    Taylor expansion is exact.
  \item \textbf{Logistic regression:} $\grad_{\bw}\Remp = \bX\T(\ba - \by)$,
    $\hess_{\bw}\Remp = \bX\T\bR\bX$, so
    \[
      \bw^{\text{new}} = \bw^{\text{old}} - (\bX\T\bR\bX)\inv\bX\T(\ba - \by).
    \]
    This is \textbf{iteratively reweighted least squares}: every step solves
    a weighted least-squares problem, and the weights $\bR$ change in each
    iteration because they depend on the current predictions.
\end{itemize}
\end{thmbox}

\begin{notebox}[title={Logistic regression on separable data}]
If the two classes are \emph{linearly separable}, the binary cross-entropy
can be pushed arbitrarily close to zero by scaling up any separating
$\bw$: the sigmoid becomes steeper and steeper, and the maximum-likelihood
solution does not exist ($\norm{\bw} \to \infty$). On non-separable data (the
typical case, shown on the slides) a finite optimum exists. Regularization
(\cref{ch:regularization}) is the standard remedy for the separable case.
\end{notebox}

\paragraph{A real-world example}
In 2019, the Austrian public employment service (AMS) announced that an
algorithm would be used to assess job-seekers. The model was a logistic
regression on several person-centric variables (gender, age, occupational
group, education, citizenship group, impairment, employment history, and
their interactions), estimating a job-seeker's short-term chance of
re-employment (90 days of employment within seven months). It was shown to be
unfair, for example with respect to women. The example is a reminder that
even the simplest models in this chapter are used for high-stakes decisions
--- and that a model fitted by maximum likelihood faithfully reproduces the
patterns, including the biases, of its training data.


% ----------------------------------------------------------
\section{Linear Neuron with Softmax Activation: Softmax Regression}
\label{sec:softmax}
% ----------------------------------------------------------

Next, the labels can take one of $K$ classes: cars, dogs, cats, airplanes,
\ldots\ in image recognition. Coding the classes as numbers $1, \dots, K$ and
running regression would impose a meaningless order (is a cat ``between''
a dog and an airplane?). Instead, each label becomes a vector.

\begin{defbox}[title={One-hot encoding}]
A label of class $k$ out of $K$ classes is coded as the vector
$\by \in \{0,1\}^K$ with $y_k = 1$ and all other entries $0$. For example,
with classes A, B, C, D: $\text{A} \mapsto (1,0,0,0)$, $\text{C} \mapsto
(0,0,1,0)$. The same encoding is used for categorical \emph{input}
features (\cref{sec:tricks}).
\end{defbox}

\begin{defbox}[title={Softmax regression}]
The network has $K$ output units and a \emph{weight matrix}
$\bW \in \R^{K \times D}$ instead of a weight vector:
\[
  g(\bx; \bW) = \ba = \softmax(\bW\bx), \qquad
  \softmax(\bs) = \Bigl(\frac{e^{s_1}}{\sum_{j=1}^K e^{s_j}}, \dots,
                        \frac{e^{s_K}}{\sum_{j=1}^K e^{s_j}}\Bigr)\T .
\]
The output is a \textbf{probability vector}: all entries are positive and
sum to one, and $a_k$ is the predicted probability of class $k$. The weight
$w_{kd}$ connects input $d$ to output $k$; row $k$ of $\bW$ is the weight
vector $\bw_k$ of output unit $k$.
\end{defbox}

\subsection{Categorical Cross-Entropy}

Same recipe once more. The probability the model assigns to a one-hot label
is $p(\by \given \ba) = \prod_{k=1}^K a_k^{y_k}$ (only the factor of the
correct class is not 1), and its negative logarithm is the loss.

\begin{defbox}[title={Categorical cross-entropy}]
\[
  \loss(\by, \ba) = -\log p(\by \given \ba) = -\sum_{k=1}^{K} y_k \log a_k,
  \qquad
  \Remp(\bW, \bX, \bY) = -\sum_{n=1}^{N}\sum_{k=1}^{K} y^n_k \log a^n_k .
\]
For a one-hot label, only one term survives: the loss is minus the log of
the predicted probability of the correct class.
\end{defbox}

\subsection{Gradient}

The derivative of the softmax is \emph{coupled}: each output depends on all
pre-activations $s_j = \bw_j\T\bx$,
\[
  \pd{a_k}{s_j} = a_k(\delta_{kj} - a_j)
  \qquad (\delta_{kj} = 1 \text{ if } k = j, \text{ else } 0).
\]
Combining with the cross-entropy and using $\sum_k y_k = 1$:
\[
  \pd{\loss}{s_j} = -\sum_{k=1}^K y_k \frac{1}{a_k}\,a_k(\delta_{kj} - a_j)
  = -y_j + a_j\sum_{k} y_k = a_j - y_j .
\]
Since $\partial s_\xi / \partial w_{\xi d} = x_d$, the derivative with
respect to a single weight and the gradient with respect to the whole
matrix are
\[
  \pd{\loss}{w_{\xi d}} = (a_\xi - y_\xi)\,x_d,
  \qquad
  \grad_{\bW}\loss = (\ba - \by)\,\bx\T .
\]
(The full derivation is part of the exercises.)

\begin{exbox}[title={Softmax and cross-entropy with numbers}]
Three classes, pre-activations $\bs = \bW\bx = (2, 1, 0.1)\T$, and the true
class is the second one, $\by = (0, 1, 0)\T$.
\begin{itemize}
  \item $e^{\bs} = (7.389, 2.718, 1.105)$, sum $11.212$, so
    $\ba = \softmax(\bs) = (0.659, 0.242, 0.099)\T$ (sums to 1).
  \item Loss: $-\log a_2 = -\log 0.242 = 1.417$.
  \item Error signal: $\ba - \by = (0.659, -0.758, 0.099)\T$. The gradient
    $(\ba - \by)\bx\T$ pushes the weights of class 2 \emph{up} (towards $\bx$)
    and those of classes 1 and 3 \emph{down}, the more so the larger their
    wrong probability.
\end{itemize}
\end{exbox}

\paragraph{Type of optimization problem and convergence}
As for logistic regression, the categorical cross-entropy is convex in
$\bW$, there is no closed-form solution, and gradient descent (or Newton's
method) is used. Two remarks answer the slides' questions about uniqueness
and convergence:
\begin{itemize}
  \item The softmax does not change if the same number is added to all
    pre-activations, $\softmax(\bs + c\ones) = \softmax(\bs)$. Adding the same
    vector $\bm{v}$ to every row of $\bW$ adds $\bm{v}\T\bx$ to every $s_k$,
    so infinitely many weight matrices give the same model: the minimizer is
    never unique (it is unique only up to such shifts).
  \item As in the binary case, perfectly separable data let the weights grow
    without bound.
\end{itemize}

\begin{intbox}
For $K = 2$ the softmax reduces to the sigmoid:
$a_1 = e^{s_1}/(e^{s_1} + e^{s_2}) = 1/(1 + e^{-(s_1 - s_2)}) = \sigma(s_1 -
s_2)$. Softmax regression is logistic regression with $K$ classes; only the
\emph{difference} of the pre-activations matters --- which is exactly the
shift invariance above.
\end{intbox}

\subsection{Comparing the Single-Layer Networks}

Put the three gradients next to each other:

\begin{thmbox}[title={Same gradient, three models}]
\begin{align*}
  \text{linear regression:}  \quad & \grad_{\bw}\Remp(y, \bx, \bw) = (\bw\T\bx - y)\,\bx, \\
  \text{logistic regression:}\quad & \grad_{\bw}\Remp(y, \bx, \bw) = (\sigma(\bw\T\bx) - y)\,\bx, \\
  \text{softmax regression:} \quad & \grad_{\bW}\Remp(\by, \bx, \bW) = (\softmax(\bW\bx) - \by)\,\bx\T .
\end{align*}
In all three cases: \textbf{(prediction $-$ target) $\times$ input}. The
pairs (linear output, squared loss), (sigmoid output, binary cross-entropy)
and (softmax output, categorical cross-entropy) are ``matched'' --- in
statistics they are called \emph{canonical links} --- and the derivative of
the activation cancels against the derivative of the loss.
\end{thmbox}

\begin{table}[htbp]
\centering
\small
\begin{tabular}{@{}lllll@{}}
\toprule
\textbf{Network} & \textbf{Labels} & \textbf{Output} & \textbf{Loss} & \textbf{Learning} \\
\midrule
linear neuron & $y \in \R$ & $\bw\T\bx$ & squared loss &
  closed form, GD, SGD \\
perceptron & $y \in \{-1, 1\}$ & $\sign(\bw\T\bx)$ & perceptron loss &
  perceptron rule \\
logistic regression & $y \in \{0, 1\}$ & $\sigma(\bw\T\bx)$ & binary CE &
  GD, SGD, Newton/IRLS \\
softmax regression & $\by \in \{0,1\}^K$ & $\softmax(\bW\bx)$ & categorical CE &
  GD, SGD, Newton \\
\bottomrule
\end{tabular}
\caption{The four single-layer networks of Sections 4.1--4.4. For many
applications these networks are already sufficient --- but they all have
\emph{linear} decision boundaries, which is their fundamental limitation
(\cref{sec:xor}).}
\label{tab:single-layer}
\end{table}

% ----------------------------------------------------------
\section{Single-Layer Networks Cannot Solve XOR}
\label{sec:xor}
% ----------------------------------------------------------

The lecture opens the multi-layer part with a quick task. Here are some
labelled points --- what are the missing labels?
\[
\begin{array}{rcl@{\qquad}rcl}
(-0.02,\ 0.01) &\to& 0 & (0.02,\ 0.98) &\to& 1 \\
(1.05,\ 0.99)  &\to& 0 & (1.03,\ -0.02) &\to& 1 \\
(0.03,\ 0.01)  &\to& 0 & (1.03,\ 1.01) &\to& 0 \\
(0.97,\ -0.01) &\to& 1 & (0.96,\ 0.03) &\to& 1 \\
(1.00,\ 0.01)  &\to& \;? & (-0.04,\ 0.03) &\to& \;?
\end{array}
\]
The points cluster around the corners of the unit square, and the label is
1 exactly when one coordinate is (close to) 1 and the other (close to) 0: it
is a noisy version of \textbf{XOR}. So the answers are 1 and 0. Easy for a
human --- and impossible for every network of the previous sections.

\subsection{Why a Linear Network Fails}

Take the four clean XOR points
$\bx^1 = (0,0)$, $\bx^2 = (1,0)$, $\bx^3 = (0,1)$, $\bx^4 = (1,1)$ with labels
$y^1 = 0$, $y^2 = 1$, $y^3 = 1$, $y^4 = 0$, and a linear (single-layer)
network $g_1(\bx; \bw) = x_1 w_1 + x_2 w_2$. Fitting the labels exactly
requires
\[
  0 = 0\,w_1 + 0\,w_2, \qquad 1 = 1\,w_1 + 0\,w_2, \qquad
  1 = 0\,w_1 + 1\,w_2, \qquad 0 = 1\,w_1 + 1\,w_2 .
\]
The second and third equation give $w_1 = w_2 = 1$, and then the fourth
says $0 = 2$. No parameters solve the problem.

The failure is more fundamental than this particular model. Adding a bias
does not help, and neither does a sigmoid or softmax on the output: these
are monotone functions of a \emph{linear} score $z(\bx) = \bw\T\bx + b$, so
the decision boundary is still a straight line. And no straight line
separates XOR. For an affine function,
$z(\bx^2) + z(\bx^3) = z(\bx^1) + z(\bx^4)$ (both sides equal
$w_1 + w_2 + 2b$). The two class-1 points would need scores above the
threshold and the two class-0 points scores below it --- impossible when
the two sums are equal.

\subsection{A Hidden Layer Solves It}

The way out is a hidden layer with a non-linearity. The lecture uses a
two-layer network with ReLU hidden units:
\[
  g_2(\bx; \bW, \bw, \bb) = \bw\T\max(0, \bW\bx + \bb) = \bw\T\bh,
  \qquad
  \bW = \begin{pmatrix} 1 & 1 \\ 1 & 1 \end{pmatrix}, \;
  \bb = \begin{pmatrix} 0 \\ -1 \end{pmatrix}, \;
  \bw = \begin{pmatrix} 1 \\ -2 \end{pmatrix}.
\]

\begin{exbox}[title={The XOR network, step by step}]
\centering
\begin{tabular}{@{}cccccc@{}}
\toprule
$\bx$ & $\bW\bx + \bb$ & $\bh = \max(0, \bW\bx + \bb)$ & $\bw\T\bh = h_1 - 2h_2$ & label \\
\midrule
$(0,0)$ & $(0, -1)$ & $(0, 0)$ & $0$ & 0 \\
$(1,0)$ & $(1, 0)$  & $(1, 0)$ & $1$ & 1 \\
$(0,1)$ & $(1, 0)$  & $(1, 0)$ & $1$ & 1 \\
$(1,1)$ & $(2, 1)$  & $(2, 1)$ & $0$ & 0 \\
\bottomrule
\end{tabular}

\raggedright
\smallskip
The network reproduces all four labels exactly. The first hidden unit
counts how many inputs are on; the second fires only if both are on, and
the output subtracts it twice.
\end{exbox}

\begin{figure}[htbp]
\centering
\begin{tikzpicture}[
    c0/.style={circle, draw=figred, fill=figred!25, thick, minimum size=2.6mm, inner sep=0pt},
    c1/.style={circle, draw=figblue, fill=figblue!30, thick, minimum size=2.6mm, inner sep=0pt}]
  % input space
  \begin{scope}
    \draw[gray!60, -{Stealth[length=1.6mm]}] (-0.4,0) -- (2.6,0) node[right, font=\scriptsize] {$x_1$};
    \draw[gray!60, -{Stealth[length=1.6mm]}] (0,-0.4) -- (0,2.6) node[above, font=\scriptsize] {$x_2$};
    \node[c0] at (0,0) {}; \node[c0] at (2,2) {};
    \node[c1] at (2,0) {}; \node[c1] at (0,2) {};
    \node[font=\tiny, below left] at (0,0) {$(0,0)$};
    \node[font=\tiny, below] at (2,-0.1) {$(1,0)$};
    \node[font=\tiny, left] at (-0.1,2) {$(0,1)$};
    \node[font=\tiny, right] at (2.1,2) {$(1,1)$};
    \node[font=\small, align=center] at (1.1,-1.0) {input space:\\no line separates the classes};
  \end{scope}
  \draw[-{Stealth[length=2.5mm]}, thick, accent] (3.3,1.1) -- node[above, font=\scriptsize] {$\bh = \max(0, \bW\bx + \bb)$} (5.9,1.1);
  % hidden space
  \begin{scope}[xshift=6.9cm]
    \draw[gray!60, -{Stealth[length=1.6mm]}] (-0.4,0) -- (4.6,0) node[right, font=\scriptsize] {$h_1$};
    \draw[gray!60, -{Stealth[length=1.6mm]}] (0,-0.4) -- (0,2.6) node[above, font=\scriptsize] {$h_2$};
    \foreach \t in {1,2} \draw[gray!50] (2*\t,0.05) -- (2*\t,-0.05) node[below, font=\tiny] {\t};
    \draw[gray!50] (0.05,2) -- (-0.05,2) node[left, font=\tiny] {1};
    \draw[accent, thick, dashed] (1,0) -- (4.6,1.8) node[below right, font=\scriptsize] {$h_1 - 2h_2 = 0.5$};
    \node[c0] at (0,0) {};
    \node[c1] at (2,0) {};
    \node[font=\tiny, above=1pt] at (2,0.12) {(two points)};
    \node[c0] at (4,2) {};
    \node[font=\small, align=center] at (2.1,-1.0) {hidden space:\\linearly separable};
  \end{scope}
\end{tikzpicture}
\caption{What the hidden layer does. In input space (left) the XOR classes
cannot be separated by a line. The hidden layer maps $(1,0)$ and $(0,1)$ to
the same point $(1,0)$ and $(1,1)$ to $(2,1)$. In this new space a single
line --- the output unit $h_1 - 2h_2$, thresholded at $0.5$ --- separates
the classes.}
\label{fig:xor}
\end{figure}

\begin{intbox}
This is the key idea of all of deep learning in its smallest form: the
hidden layer \emph{transforms the data into a new representation} in which
the problem becomes easy (here: linearly separable), and the output layer is
just a linear model on top. In a deep network, every layer does this once
more.
\end{intbox}

\begin{thmbox}[title={Universal function approximation (Hornik, 1989; Cybenko, 1989)}]
A multi-layer perceptron with a single hidden layer, enough hidden units and
a suitable non-linear activation function (for example a sigmoid) can
approximate any continuous function on a compact set to arbitrary accuracy.
\end{thmbox}

So single-layer networks cannot solve XOR, a simple two-layer network with a
ReLU non-linearity does, and MLPs can even approximate any continuous
function. The theorem says that a solution \emph{exists}; it says nothing
about how many hidden units are needed or how to find the weights. Finding
them is the job of backpropagation.


% ----------------------------------------------------------
\section{Multi-Layer Perceptron}
\label{sec:mlp}
% ----------------------------------------------------------

\subsection{Model and Forward Pass}

A multi-layer perceptron (MLP) consists of an input layer $L_1$, one or more
hidden layers, and an output layer; every unit of one layer is connected to
every unit of the next. The lecture numbers the units consecutively through
the network (\cref{sec:notation-units}): the $D$ input units come first,
the $K$ output units are the last ones, $Q - K + 1, \dots, Q$.

\begin{defbox}[title={Quantities of an MLP}]
For every non-input unit $i$ (i.e.\ $i > D$):
\[
  s_i = \sum_{j} w_{ij}\, a_j + b_i \qquad\text{(network input, pre-activation)},
  \qquad
  a_i = f(s_i) \qquad\text{(activation)},
\]
where the sum runs over all units $j$ with a connection to $i$, and for
input units $a_i = x_i$. The \textbf{architecture} of a network is given by
the number of layers, the number of units per layer and the connections
between them; the activation functions may also be counted as part of it.
Different units can have different activation functions $f_i$.
\end{defbox}

\begin{figure}[htbp]
\centering
\begin{tikzpicture}
  \node[inunit] (x1) at (-0.8,0) {};
  \node[inunit] (x2) at (0.8,0) {$j$};
  \foreach \i/\xx in {1/-2.4,2/-0.8,3/0.8} \node[unit] (h\i) at (\xx,2) {};
  \node[unit, fill=figblue!20] (h4) at (2.4,2) {$i$};
  \node[outunit] (o) at (0,4) {};
  \foreach \i in {1,...,4} {
    \draw[conn] (x1) -- (h\i); \draw[conn] (x2) -- (h\i); \draw[conn] (h\i) -- (o);
  }
  \draw[-{Stealth[length=2.6mm]}, very thick, figblue] (x2) -- node[right, font=\small] {$w_{ij}$} (h4);
  \node[font=\small, figblue, anchor=west] at (3.0,2.25) {$s_i = \sum_j w_{ij}a_j + b_i$};
  \node[font=\small, figblue, anchor=west] at (3.0,1.65) {$a_i = f(s_i)$};
  \node[font=\small, anchor=east] at (-3.2,0) {$L_1$: input};
  \node[font=\small, anchor=east] at (-3.2,2) {$L_2$: hidden};
  \node[font=\small, anchor=east] at (-3.2,4) {$L_3$: output};
\end{tikzpicture}
\caption{An MLP with one hidden layer. Unit $i$ collects the weighted
activations of the units below (pre-activation $s_i$, including its bias
$b_i$) and passes them through its activation function.}
\label{fig:mlp}
\end{figure}

\paragraph{Activation functions}
The hidden units need a \emph{non-linear} activation function --- otherwise
the whole network collapses to a linear model (\cref{sec:dnn}). The lecture
shows six common choices; \cref{tab:activations} lists them with their
derivatives (which backpropagation needs), and \cref{fig:activations}
plots them.

\begin{table}[htbp]
\centering
\small
\renewcommand{\arraystretch}{1.25}
\begin{tabular}{@{}lll@{}}
\toprule
\textbf{Name} & $f(x)$ & $f'(x)$ \\
\midrule
logistic sigmoid & $\sigma(x) = \dfrac{1}{1+e^{-x}}$ & $\sigma(x)(1-\sigma(x)) \le \tfrac14$ \\
tanh & $\tanh(x) = \dfrac{e^x - e^{-x}}{e^x + e^{-x}} = 2\sigma(2x) - 1$ & $1 - \tanh^2(x) \le 1$ \\
ReLU & $\max(0, x)$ & $1$ for $x > 0$, $0$ for $x < 0$ \\
leaky ReLU & $x$ for $x > 0$, $\alpha x$ for $x \le 0$ (small $\alpha$) & $1$ for $x > 0$, $\alpha$ for $x < 0$ \\
ELU & $x$ for $x > 0$, $\alpha(e^x - 1)$ for $x \le 0$ & $1$ for $x > 0$, $\alpha e^x$ for $x < 0$ \\
SELU & $\lambda x$ for $x > 0$, $\lambda\alpha(e^x - 1)$ for $x \le 0$ & $\lambda$ for $x > 0$, $\lambda\alpha e^x$ for $x < 0$ \\
\bottomrule
\end{tabular}
\caption{Common activation functions. For SELU, $\lambda \approx 1.0507$ and
$\alpha \approx 1.6733$ are fixed constants (chosen so that activations stay
normalized, see \cref{ch:normalization}).}
\label{tab:activations}
\end{table}

\begin{figure}[htbp]
\centering
\begin{tikzpicture}
\begin{groupplot}[
    group style={group size=3 by 2, horizontal sep=9mm, vertical sep=11mm},
    dlbt, width=0.36\linewidth, height=3.7cm,
    xmin=-3, xmax=3, ymin=-1.9, ymax=3.1,
    xtick={-2,0,2}, ytick={-1,0,1,2,3},
    title style={font=\small, yshift=-1.5mm}]
  \nextgroupplot[title={sigmoid}]
    \addplot[very thick, accent, domain=-3:3] {1/(1+exp(-x))};
    \addplot[thick, figamber, dashed, domain=-3:3] {exp(-x)/(1+exp(-x))^2};
  \nextgroupplot[title={tanh}]
    \addplot[very thick, accent, domain=-3:3] {tanh(x)};
    \addplot[thick, figamber, dashed, domain=-3:3] {1-tanh(x)^2};
  \nextgroupplot[title={ReLU}]
    \addplot[very thick, accent, domain=-3:3] {max(0,x)};
    \addplot[thick, figamber, dashed, domain=-3:0] {0};
    \addplot[thick, figamber, dashed, domain=0:3] {1};
  \nextgroupplot[title={leaky ReLU ($\alpha = 0.1$)}]
    \addplot[very thick, accent, domain=-3:3] {max(0.1*x,x)};
    \addplot[thick, figamber, dashed, domain=-3:0] {0.1};
    \addplot[thick, figamber, dashed, domain=0:3] {1};
  \nextgroupplot[title={ELU ($\alpha = 1$)}]
    \addplot[very thick, accent, domain=-3:0] {exp(x)-1};
    \addplot[very thick, accent, domain=0:3] {x};
    \addplot[thick, figamber, dashed, domain=-3:0] {exp(x)};
    \addplot[thick, figamber, dashed, domain=0:3] {1};
  \nextgroupplot[title={SELU}]
    \addplot[very thick, accent, domain=-3:0] {1.0507*1.6733*(exp(x)-1)};
    \addplot[very thick, accent, domain=0:3] {1.0507*x};
    \addplot[thick, figamber, dashed, domain=-3:0] {1.0507*1.6733*exp(x)};
    \addplot[thick, figamber, dashed, domain=0:3] {1.0507};
\end{groupplot}
\end{tikzpicture}
\caption{Activation functions (solid) and their derivatives (dashed). The
sigmoid's derivative never exceeds $0.25$ and both sigmoid and tanh saturate
for large $|x|$. ReLU and its variants have derivative $1$ (SELU:
$\lambda$) for all positive inputs --- the property that makes them
attractive for deep networks (\cref{sec:vgp}).}
\label{fig:activations}
\end{figure}

\paragraph{Output activations}
The activation of the \emph{output} units is chosen according to the task,
exactly as in the single-layer networks: usually a \textbf{linear}
activation for regression, the \textbf{logistic sigmoid} for binary (or
multi-label) classification, and the \textbf{softmax} for multi-class
classification. An MLP is therefore ``a linear, logistic or softmax
regression on top of learned hidden features''.

\paragraph{Forward pass}
Computing the output for an input is called the \emph{forward pass}: start
with the input units and go layer by layer upwards, computing $s_i$ and
$a_i$ for every unit.

\begin{algorithm}[htbp]
\caption{Forward pass of an MLP}
\label{alg:forward}
\KwIn{input $\bx$, weights $w_{ij}$, layers $L_1, \dots, L_L$}
\For(\tcp*[f]{initialization}){$i = 1, \dots, D$}{
  $a_i \leftarrow x_i$\;
}
\For{$\nu = 2, \dots, L$}{
  \For{all $i \in L_\nu$}{
    $s_i \leftarrow \sum_{j:\, w_{ij}\text{ exists}} w_{ij}\, a_j$\;
    $a_i \leftarrow f(s_i)$\;
  }
}
\KwOut{$\hat y_i = (g(\bx; \bw))_i = a_i$ for all $Q - K + 1 \le i \le Q$}
\end{algorithm}

\subsection{Gradient Descent and the Backpropagation Algorithm}
\label{sec:backprop}

To train an MLP we minimize the empirical error by gradient descent, just as
before:
\[
  \Remp(\bw, \bX, \bY) = \frac1N\sum_{n=1}^{N} \loss\bigl(\by^n, g(\bx^n; \bw)\bigr),
  \qquad
  \bw^{\text{new}} = \bw^{\text{old}} - \eta\,\grad_{\bw}\Remp,
\]
where $g(\bx^n; \bw)$ is the network and $\bw$ contains all weights (weight
matrices) of all layers and the biases. For a single data point the update
uses $\grad_{\bw}\loss(\by, g(\bx; \bw))$.

The only new problem is computing the gradient: a weight deep inside the
network influences the loss through all the layers above it.
\textbf{Backpropagation} --- made popular by Rumelhart et al.\ (1986),
proposed earlier by Werbos (1974), also called \emph{delta propagation} ---
solves it with nothing but the chain rule, organized cleverly.

\paragraph{Step 1: split the derivative at the pre-activation}
Weight $w_{ij}$ influences the loss only through the pre-activation $s_i =
\sum_j w_{ij}a_j + b_i$ of the unit it points to. By the chain rule,
\[
  \pd{\loss}{w_{ij}} = \pd{\loss}{s_i}\,\pd{s_i}{w_{ij}} = \pd{\loss}{s_i}\, a_j,
  \qquad\text{since } \pd{s_i}{w_{ij}} = a_j .
\]

\begin{defbox}[title={Delta error}]
The \textbf{delta error} of unit $i$ is the derivative of the loss with
respect to its pre-activation:
\[
  \delta_i \coloneqq \pd{}{s_i}\,\loss\bigl(\by^n, g(\bx^n; \bw)\bigr).
\]
With it, the gradient for every weight is simply
\[
  \pd{\loss}{w_{ij}} = \delta_i\, a_j
  \qquad\text{(delta at the target unit $\times$ activation at the source unit).}
\]
\end{defbox}

Compare this with softmax regression: there,
$\partial\Remp / \partial w_{ij} = (\softmax(\bW\bx)_i - y_i)\,x_j$, i.e.\
$\delta_i = a_i - y_i$ and $a_j = x_j$. The single-layer result is the
special case of a network without hidden units.

\paragraph{Step 2: deltas at the output units}
For an output unit $k$ (with an element-wise activation),
\[
  \delta_k = \pd{\loss}{a_k}\, f'(s_k),
  \qquad Q - K + 1 \le k \le Q,
\]
where $a_k = (g(\bx^n;\bw))_k$ is the output activation and $s_k$ the
output pre-activation. For the matched loss/activation pairs this is the
familiar error signal (\cref{tab:loss-activation}) --- the same as for
linear, logistic and softmax regression.

\paragraph{Step 3: deltas at the hidden units (the recursion)}
A hidden unit $j$ influences the loss only through the units $i$ in the
layer above that it is connected to, via
$s_i = \sum_j w_{ij} f(s_j) + b_i$. The multivariate chain rule therefore
gives
\[
  \delta_j = \pd{\loss}{s_j} = \sum_i \pd{\loss}{s_i}\pd{s_i}{s_j}
  = \sum_i \delta_i\, w_{ij} f'(s_j)
  = f'(s_j)\sum_i \delta_i\, w_{ij},
\]
where the sum goes over all units $i$ for which the weight $w_{ij}$ exists.

\begin{thmbox}[title={Backpropagation (unit-wise)}]
\[
  \underbrace{\delta_k = \pd{\loss}{a_k}\,f'(s_k)}_{\text{output units}},
  \qquad
  \underbrace{\delta_j = f'(s_j)\sum_i \delta_i\, w_{ij}}_{\text{hidden units}},
  \qquad
  \underbrace{\pd{\loss}{w_{ij}} = \delta_i\, a_j, \quad
  \Delta w_{ij} = -\eta\,\delta_i\, a_j}_{\text{gradient and update}} .
\]
\end{thmbox}

\begin{figure}[htbp]
\centering
\begin{tikzpicture}[xscale=1.35, yscale=1.1]
  \foreach \i in {1,2} \node[inunit] (a\i) at (\i+0.5,0) {};
  \foreach \i in {1,...,4} \node[unit] (b\i) at (\i,1.6) {};
  \foreach \i in {1,...,4} \node[unit, fill=figred!20, draw=figred] (c\i) at (\i,3.2) {};
  \node[unit, fill=figred!35, draw=figred] (d) at (2.5,4.8) {};
  \foreach \i in {1,2} \foreach \j in {1,...,4} \draw[conn, gray!60] (a\i) -- (b\j);
  \foreach \i in {1,...,4} \foreach \j in {1,...,4} \draw[conn, gray!60] (b\i) -- (c\j);
  \foreach \i in {1,...,4} \draw[conn, gray!60] (c\i) -- (d);
  % highlight unit i in hidden 1
  \node[unit, fill=figred!12, draw=figred, very thick] (bi) at (4,1.6) {$i$};
  \foreach \j in {1,...,4} \draw[bconn] (c\j) -- (bi);
  \node[font=\scriptsize, figred, anchor=west, align=left] at (4.45,1.6)
    {$\delta_i = f'(s_i)\sum_k \delta_k w_{ki}$};
  \node[font=\scriptsize, figred, anchor=west] at (4.45,3.2) {$\delta_k$};
  \node[font=\scriptsize, figred, anchor=west] at (2.9,4.8)
    {$\delta = \frac{\partial\loss}{\partial a}\,f'(s)$};
  \node[font=\small, anchor=east] at (0.3,0) {$L_1$: input};
  \node[font=\small, anchor=east] at (0.3,1.6) {$L_2$: hidden 1};
  \node[font=\small, anchor=east] at (0.3,3.2) {$L_3$: hidden 2};
  \node[font=\small, anchor=east] at (0.3,4.8) {$L_4$: output};
\end{tikzpicture}
\caption{Delta propagation. The deltas are computed at the output and then
passed \emph{backwards} through the same weights that carried the
activations forward: unit $i$ in hidden layer 1 collects the deltas of all
units $k$ it feeds into, weighted by $w_{ki}$, and multiplies by its own
$f'(s_i)$. The gradient for a weight $w_{ki}$ is then $\delta_k a_i$.}
\label{fig:delta-propagation}
\end{figure}

\begin{table}[htbp]
\centering
\small
\renewcommand{\arraystretch}{1.25}
\begin{tabular}{@{}llll@{}}
\toprule
\textbf{Loss function} & & \textbf{Output activation} & \textbf{Delta error} \\
\midrule
squared loss & $\tfrac12(a - y)^2$ & linear: $f(x) = x$ & $\delta = a - y$ \\
binary CE & $-y\log a - (1-y)\log(1-a)$ & sigmoid: $f(x) = \frac{1}{1+e^{-x}}$ & $\delta = a - y$ \\
categorical CE & $-\sum_{k=1}^K y_k\log a_k$ & softmax: $f(\bx) = \softmax(\bx)$ & $\bdelta = \ba - \by$ \\
\bottomrule
\end{tabular}
\caption{Frequently used loss functions and their matching output
activations. For these pairs, the derivative of the activation cancels and
the delta error at the output is simply ``prediction minus target''.}
\label{tab:loss-activation}
\end{table}

\begin{notebox}[title={Potential problems of other combinations}]
The slides ask what goes wrong if loss and output activation do not match.
Two examples:
\begin{itemize}
  \item \textbf{Squared loss with a sigmoid output:} now
    $\delta = (a - y)\,\sigma'(s) = (a - y)\,a(1-a)$. If the network is
    \emph{confidently wrong} (say $a \approx 0$ but $y = 1$), then
    $a(1-a) \approx 0$ and the delta almost vanishes --- the network learns
    slowest exactly when it is most wrong. With cross-entropy the factor
    $a(1-a)$ cancels.
  \item \textbf{Cross-entropy with a linear output:} the output is not
    restricted to $(0, 1)$, so $\log a$ or $\log(1-a)$ may be undefined.
\end{itemize}
Also note that for the softmax the derivative is not element-wise:
``$\partial\loss/\partial a_k\,f'(s_k)$'' has to be replaced by the full
Jacobian of the softmax, which yields $\bdelta = \ba - \by$ in combination
with the cross-entropy (\cref{sec:softmax}).
\end{notebox}

\begin{algorithm}[htbp]
\caption{Backward pass of an MLP (online version)}
\label{alg:backward}
\KwIn{activations $a_i$ and pre-activations $s_i$ from the forward pass,
  label $\by$, learning rate $\eta$}
\For(\tcp*[f]{output layer}){$i = Q - K + 1, \dots, Q$}{
  $\delta_i \leftarrow \pd{\loss}{a_i}\, f'(s_i)$\;
  \For{all $j \in L_{L-1}$}{
    $\Delta w_{ij} \leftarrow -\eta\,\delta_i\, a_j$\;
  }
}
\For(\tcp*[f]{hidden layers, top to bottom}){$\nu = L - 1, \dots, 2$}{
  \For{all $i \in L_\nu$}{
    $\delta_i \leftarrow f'(s_i)\sum_k \delta_k\, w_{ki}$\;
    \For{all $j \in L_{\nu-1}$}{
      $\Delta w_{ij} \leftarrow -\eta\,\delta_i\, a_j$\;
    }
  }
}
\end{algorithm}

\begin{thmbox}[title={Backpropagation in five steps (recap)}]
\begin{enumerate}
  \item Select a sample $\bx$ (randomly) from the training set and perform a
    forward pass. Memorize the activations $a_i$ and pre-activations $s_i$
    of all units.
  \item Compute the delta errors at the output units:
    $\delta_k = \frac{\partial\loss}{\partial a_k}f'(s_k)$.
  \item Compute the delta errors of the hidden layers, starting with the
    layer closest to the output, by backpropagating the error signal:
    $\delta_j = f'(s_j)\sum_i \delta_i w_{ij}$.
  \item Compute the gradients $\frac{\partial\loss}{\partial w_{ij}} = \delta_i a_j$.
  \item Update the weights with a gradient-descent step:
    $\Delta w_{ij} = -\eta\,\delta_i\, a_j$.
\end{enumerate}
\end{thmbox}

\begin{exbox}[title={Backpropagation with numbers (checked by finite differences)}]
A 2--2--1 network: sigmoid hidden units, a linear output unit, squared loss,
no biases. Input $\bx = (1, 0.5)\T$, label $y = 1$,
\[
  \bW\lay{1} = \begin{pmatrix} 0.5 & -0.3 \\ 0.2 & 0.8 \end{pmatrix}, \qquad
  \bW\lay{2} = \begin{pmatrix} 0.7 & -0.4 \end{pmatrix}.
\]
\textbf{Forward pass.}
$\bs\lay{1} = \bW\lay{1}\bx = (0.5 - 0.15,\; 0.2 + 0.4)\T = (0.35, 0.60)\T$,
$\ba\lay{1} = \sigma(\bs\lay{1}) = (0.5866, 0.6457)\T$,
$\hat y = 0.7 \cdot 0.5866 - 0.4 \cdot 0.6457 = 0.1524$,
loss $\loss = \frac12(0.1524 - 1)^2 = 0.3592$.

\textbf{Output delta.} Linear output with squared loss:
$\delta\lay{2} = \hat y - y = -0.8476$.

\textbf{Hidden deltas.} $f'(\bs\lay{1}) = \ba\lay{1}\odot(1 - \ba\lay{1})
= (0.2425, 0.2288)\T$, so
\[
  \delta\lay{1}_1 = 0.2425 \cdot 0.7 \cdot (-0.8476) = -0.1439, \qquad
  \delta\lay{1}_2 = 0.2288 \cdot (-0.4) \cdot (-0.8476) = 0.0776 .
\]
\textbf{Gradients} ($\delta$ at the target times activation at the source):
\[
  \pd{\loss}{\bW\lay{2}} = \delta\lay{2}\,\ba\lay{1\mathsf{T}} = (-0.4972,\; -0.5473),
  \qquad
  \pd{\loss}{\bW\lay{1}} = \bdelta\lay{1}\bx\T =
  \begin{pmatrix} -0.1439 & -0.0719 \\ 0.0776 & 0.0388 \end{pmatrix}.
\]
Central finite differences give the same six numbers (maximum deviation
below $10^{-10}$). One update with $\eta = 0.5$ raises the output to
$\hat y = 0.4968$ and lowers the loss from $0.359$ to $0.127$.
\end{exbox}

\paragraph{Why backprop is efficient}
The name refers to the fact that the delta errors of one layer are used to
compute the deltas of the layer below: every delta is computed once and
reused by all weights below it. For a network with $W$ weights, a forward
and a backward pass each cost $\bigO{W}$ operations, so the \emph{whole}
gradient costs $\bigO{W}$. Computing it naively by finite differences would
need a separate forward pass for every single weight, i.e.\ $\bigO{W^2}$. A
similar approach can be used to compute second derivatives.

\paragraph{Full batch, online and mini-batch learning}
The gradient can be computed on the full training set before taking one
step (\emph{full-batch} gradient), on a single sample followed by an update
(\emph{online learning}), or on a small random subset of, e.g., 32 samples
(\emph{mini-batch training} or \emph{stochastic gradient descent}).
Details in \cref{sec:tricks}.

\subsection{Training MLPs}

With MLPs we suddenly have many design choices: the number of layers, the
number of neurons per layer, the activation function, the learning rate.
These are \textbf{hyperparameters} --- they are not learned by gradient
descent. And minimizing the empirical error is not sufficient: what we care
about is the error on \emph{new} data (\cref{ch:generalization}). So the
architecture and hyperparameters are optimized on a separate
\emph{validation set}, and the final model is evaluated on a \emph{test set}
(\cref{sec:tricks}).

\paragraph{The loss surface}
For linear and logistic regression the error was convex. For an MLP the
loss surface as a function of the weights is \textbf{highly non-convex},
with many local minima and saddle points. The minimum is also never unique,
because of \textbf{weight-space symmetries}:
\begin{itemize}
  \item For an odd activation such as $\tanh$ ($\tanh(-x) = -\tanh(x)$),
    flipping the signs of all weights into and out of a hidden unit gives
    exactly the same network function.
  \item The order of the hidden units within a layer can be permuted (along
    with their weights) without changing the function.
\end{itemize}
A layer with $M$ tanh units therefore has $M!\,2^M$ equivalent weight
settings for every solution. Despite all this, there are theoretical results
indicating that SGD still finds ``good'' minima (more in the follow-up
course).


% ----------------------------------------------------------
\section{(Deep) Feed-Forward Neural Networks}
\label{sec:dnn}
% ----------------------------------------------------------

Deep feed-forward networks (DNNs) work on the same principle as MLPs ---
layers of interconnected neurons --- but with more layers, more neurons, and
different activation functions. The boundary between the two is informal
(\cref{tab:mlp-dnn}).

\begin{table}[htbp]
\centering
\small
\begin{tabular}{@{}lll@{}}
\toprule
 & \textbf{MLP} & \textbf{DNN} \\
\midrule
number of hidden layers       & 1 or few        & $> 1$, up to several hundreds \\
number of neurons per layer   & $< 10$          & $> 100$ \\
activation function           & tanh, sigmoid   & ReLU, SELU \\
\bottomrule
\end{tabular}
\caption{Typical MLPs versus deep feed-forward networks (informal
guidelines from the lecture).}
\label{tab:mlp-dnn}
\end{table}

For deep networks the unit-by-unit numbering becomes unwieldy, so we switch
to the \textbf{layer-wise notation} of \cref{sec:notation-layers}: the input
is $\ba\lay{0} = \bx$, layer $l$ has weight matrix $\bW\lay{l}$,
pre-activations $\bs\lay{l}$ and activations $\ba\lay{l}$, and the
activation function $f$ is applied element-wise. (In the MLP picture the
units are numbered $1, \dots, Q$ snaking through the network; in the DNN
picture each layer is numbered separately.)

\subsection{Forward Pass}

From one layer to the next:
\[
  \ba\lay{l} = f(\bs\lay{l}) = f(\bW\lay{l}\ba\lay{l-1})
  \qquad\text{or, with bias units,}\qquad
  \ba\lay{l} = f(\bW\lay{l}\ba\lay{l-1} + \bb\lay{l}).
\]
The whole network is a nested function,
\[
  \byhat = g(\bx; \bW\lay{1}, \dots, \bW\lay{L})
  = \sigma\Bigl(\bW\lay{L}\bigl(\cdots f\bigl(\bW\lay{2} f(\bW\lay{1}\bx)\bigr)\cdots\bigr)\Bigr),
\]
with the output activation $\sigma$ (linear, sigmoid or softmax) chosen
for the task.

\paragraph{Why non-linearities?}
Without the non-linear $f$ in the hidden layers, the nested function
collapses:
\[
  g(\bx; \bW\lay{1}, \dots, \bW\lay{L})
  = \sigma\bigl(\underbrace{\bW\lay{L}\cdots\bW\lay{1}}_{=\,\hat\bW}\,\bx\bigr)
  = \sigma(\hat\bW\bx).
\]
A product of matrices is just another matrix, so a deep \emph{linear}
network is no more powerful than a single layer --- it cannot even solve
XOR. Depth only pays off with non-linear activations.

\subsection{Backpropagation in DNNs}
\label{sec:backprop-dnn}

Without loss of generality, take a weight in the \emph{first} layer, the one
furthest from the loss. The chain rule through all layers gives a product of
Jacobians:
\[
  \pd{}{w\lay{1}_{ij}}\,\loss\bigl(\by, g(\bx; \bW\lay{1}, \dots, \bW\lay{L})\bigr)
  = \underbrace{\pd{\loss(\by, \byhat)}{\byhat}\,\pd{\byhat}{\bs\lay{L}}}_{=:\,A}
    \cdots
    \underbrace{\pd{\bs\lay{l}}{\bs\lay{l-1}}}_{=:\,B}
    \cdots
    \underbrace{\pd{\bs\lay{1}}{w\lay{1}_{ij}}}_{=:\,C} .
\]
(Derivatives of scalars with respect to vectors are row vectors here, so
the product is read left to right.) We treat the three kinds of factors in
turn, after defining the layer-wise delta errors as the transposed
derivative of the loss with respect to the pre-activations:
\[
  \bdelta\lay{l\mathsf{T}} \coloneqq \pd{\loss(\by, \byhat)}{\bs\lay{l}} .
\]

\paragraph{A --- the output layer}
$A$ depends on the choice of loss and output activation. For the matched
pairs of \cref{tab:loss-activation} (canonical links) it is simply
\[
  \bdelta\lay{L\mathsf{T}} = A = (\byhat - \by)\T .
\]

\paragraph{B --- from one layer to the one below}
Since $\bs\lay{l} = \bW\lay{l} f(\bs\lay{l-1})$, the Jacobian of one layer's
pre-activations with respect to the previous layer's is
\[
  \pd{\bs\lay{l}}{\bs\lay{l-1}} = \bW\lay{l}\,\diag\bigl(f'(\bs\lay{l-1})\bigr)
  \eqdef \bJ\lay{l}.
\]
This gives the recursive formula for the deltas:
\[
  \bdelta\lay{l-1\mathsf{T}} = \pd{\loss}{\bs\lay{l}}\pd{\bs\lay{l}}{\bs\lay{l-1}}
  = \bdelta\lay{l\mathsf{T}}\,\bW\lay{l}\diag\bigl(f'(\bs\lay{l-1})\bigr)
  = \bdelta\lay{l\mathsf{T}}\bJ\lay{l},
\]
or, transposed and written element-wise,
\[
  \bdelta\lay{l-1} = f'(\bs\lay{l-1}) \odot \bigl(\bW\lay{l\mathsf{T}}\bdelta\lay{l}\bigr).
\]
The second form is the matrix version of
$\delta_j = f'(s_j)\sum_i\delta_i w_{ij}$ from the MLP.

\paragraph{C --- the weight itself}
$\bs\lay{1} = \bW\lay{1}\ba\lay{0}$, and $w\lay{1}_{ij}$ only appears in
component $i$:
\[
  \pd{\bs\lay{1}}{w\lay{1}_{ij}} = (0, \dots, 0, a\lay{0}_j, 0, \dots, 0)\T
  \quad\text{(non-zero entry at position } i\text{)}.
\]

\paragraph{Putting it together}
For any layer $l$,
\[
  \pd{\loss}{w\lay{l}_{ij}} = \pd{\loss}{\bs\lay{l}}\pd{\bs\lay{l}}{w\lay{l}_{ij}}
  = \bigl(\bdelta\lay{l}\bigr)_i\,\bigl(\ba\lay{l-1}\bigr)_j,
  \qquad\text{or as an outer product:}\qquad
  \pd{\loss}{\bW\lay{l}} = \bdelta\lay{l}\,\ba\lay{l-1\mathsf{T}} .
\]

\begin{thmbox}[title={Backpropagation in deep networks}]
\begin{enumerate}
  \item Compute the deltas at the output: $\bdelta\lay{L}$ (e.g.\ $\byhat - \by$).
  \item Backpropagate them through the network:
    \[
      \bdelta\lay{l-1\mathsf{T}} = \bdelta\lay{l\mathsf{T}}\,
      \underbrace{\bW\lay{l}\diag\bigl(f'(\bs\lay{l-1})\bigr)}_{\bJ\lay{l}} .
    \]
  \item Compute the weight updates:
    $\Delta\bW\lay{l} = -\eta\,\bdelta\lay{l}\,\ba\lay{l-1\mathsf{T}}$.
\end{enumerate}
\end{thmbox}

\begin{notebox}[title={Shape check: $\bdelta\,\ba\T$, not $\ba\,\bdelta\T$}]
The gradient with respect to $\bW\lay{l} \in \R^{n_l \times n_{l-1}}$ must
have the same shape as $\bW\lay{l}$. Since $\bdelta\lay{l} \in \R^{n_l}$ and
$\ba\lay{l-1} \in \R^{n_{l-1}}$, the outer product is
$\bdelta\lay{l}\ba\lay{l-1\mathsf{T}}$ --- entry $(i,j)$ is
$\delta_i a_j$. One slide writes $\ba\lay{l-1}\bdelta\lay{l\mathsf{T}}$, which
is the transpose; the update rule on the following slide uses the correct
$\bdelta\lay{l}(\ba\lay{l-1})\T$. When in doubt, check the shape.
\end{notebox}

\subsection{The Vanishing Gradient Problem}
\label{sec:vgp}

The recursion $\bdelta\lay{l-1\mathsf{T}} = \bdelta\lay{l\mathsf{T}}\bJ\lay{l}$
reveals the central problem of deep networks. Unrolled from the output down
to the first layer,
\[
  \bdelta\lay{1\mathsf{T}} = \bdelta\lay{L\mathsf{T}}\,\bJ\lay{L}\bJ\lay{L-1}\cdots\bJ\lay{2},
\]
the delta errors are multiplied by one Jacobian per layer. If the norms of
all these Jacobians are bounded by some $k$, $\norm{\bJ\lay{l}} \le k$, then
\[
  \norm{\bdelta\lay{1}} \le k^{L-1}\,\norm{\bdelta\lay{L}} .
\]
The size of the error signal changes \emph{exponentially} with depth. For
$k < 1$ it vanishes: the lower layers receive almost no gradient and do not
learn.

\begin{thmbox}[title={Vanishing and exploding gradients (Hochreiter, 1991)}]
Because the deltas are multiplied by the Jacobian $\bJ\lay{l} =
\bW\lay{l}\diag(f'(\bs\lay{l-1}))$ in every layer, their norm grows or
shrinks exponentially across layers. Three regimes:
\begin{itemize}
  \item $\norm{\bdelta\lay{l-1}} < \norm{\bdelta\lay{l}}$: \textbf{vanishing
    gradients.} The typical case for sigmoid networks: since
    $|\sigma'(x)| \le 0.25$, $\norm{\diag(f'(\bs\lay{l-1}))} \le 0.25$, and
    the deltas shrink in every layer unless the weights compensate.
  \item $\norm{\bdelta\lay{l-1}} \approx \norm{\bdelta\lay{l}}$:
    \textbf{stable gradients.} The ideal case; many algorithmic improvements
    (e.g.\ initialization strategies, \cref{ch:initialization}) aim at
    keeping this ratio close to one.
  \item $\norm{\bdelta\lay{l-1}} > \norm{\bdelta\lay{l}}$: \textbf{exploding
    gradients.} If the weight matrices have large norm, the deltas grow
    through the layers; huge weight updates make learning unstable or cause
    numerical overflows.
\end{itemize}
\end{thmbox}

\begin{figure}[htbp]
\centering
\begin{tikzpicture}
\begin{axis}[axis lines=left, width=0.72\linewidth, height=5.4cm,
    ymode=log, xmin=0, xmax=20, ymin=1e-12, ymax=30,
    xlabel={number of layers $l$ the delta passes}, ylabel={upper bound $k^{l}$},
    label style={font=\small},
    tick label style={font=\scriptsize, gray!70!black},
    grid=major, grid style={gray!15},
    axis line style={gray!70},
    ytick={1e-12,1e-9,1e-6,1e-3,1},
    legend style={font=\scriptsize, draw=none, at={(0.5,-0.32)}, anchor=north},
    legend columns=2, legend cell align=left,
    samples at={0,1,...,20}]
  \addplot[very thick, figred, mark=*, mark size=1pt] {0.25^x};
  \addlegendentry{$k = 0.25$ (sigmoid, $\norm{\bW} \le 1$)\quad}
  \addplot[very thick, figamber, mark=*, mark size=1pt] {0.9^x};
  \addlegendentry{$k = 0.9$}
  \addplot[very thick, figgreen, mark=*, mark size=1pt] {1^x};
  \addlegendentry{$k = 1$ (stable)\quad}
  \addplot[very thick, figpurple, mark=*, mark size=1pt] {1.1^x};
  \addlegendentry{$k = 1.1$ (exploding)}
\end{axis}
\end{tikzpicture}
\caption{How much of the error signal survives $l$ layers, for different
per-layer factors $k$ (log scale). With the sigmoid's factor $0.25$, ten
layers already shrink the delta by a factor of about $4\cdot 10^{-6}$.}
\label{fig:vgp}
\end{figure}

\begin{exbox}[title={How fast do gradients vanish?}]
With sigmoid units and weight matrices of norm at most one, $k = 0.25$. The
delta at the first layer of a network with $L$ layers is then at most
$0.25^{L-1}$ times the output delta:
\[
  L = 5:\ 0.25^4 \approx 3.9\cdot 10^{-3}, \quad
  L = 10:\ 0.25^9 \approx 3.8\cdot 10^{-6}, \quad
  L = 20:\ 0.25^{19} \approx 3.6\cdot 10^{-12}.
\]
In a 20-layer sigmoid network, the first layers would learn about a
trillion times more slowly than the last one.
\end{exbox}

\paragraph{Ameliorations for feed-forward networks (outlook)}
The lecture names several remedies, each covered later:
\begin{itemize}
  \item \textbf{rectified linear units} --- $f'(x) = 1$ for all positive
    inputs, so the activation function does not shrink the delta;
  \item \textbf{initialization strategies} that keep the norm of the
    Jacobians close to one (\cref{sec:tricks,ch:initialization});
  \item \textbf{skip connections} (ResNets), which give the gradient a
    shortcut around layers;
  \item and more (e.g.\ normalization, \cref{ch:normalization}). For
    recurrent networks, LSTM is the classic solution (\cref{ch:history}).
\end{itemize}

% ----------------------------------------------------------
\section{Jacobian}
\label{sec:jacobian}
% ----------------------------------------------------------

``The Jacobian of a neural network'' is an ambiguous term. Depending on what
is differentiated with respect to what, it means one of three things:

\begin{defbox}[title={Three Jacobians of a neural network}]
\begin{enumerate}
  \item \textbf{Parameter Jacobian} --- the derivative of the loss with
    respect to the parameters,
    $\displaystyle \pd{\loss(\by, g(\bx; \bw))}{\bw}$.
    This is the gradient used to train the network by gradient descent
    (computed by backpropagation).
  \item \textbf{Input--output Jacobian} --- the derivative of the network
    outputs with respect to the inputs,
    $\displaystyle \pd{g(\bx; \bw)}{\bx} = \pd{\byhat}{\bx}$.
    It tells how sensitive the prediction is to each input feature and can
    be used to optimize the \emph{input} with respect to the output, by
    gradient \emph{ascent} on $\bx$ with fixed weights. ``Deep Dream''
    images are made this way: an image is changed so that it excites
    certain units more and more (a horizon turns into towers and pagodas,
    leaves into birds and insects).
  \item \textbf{(Single-layer) Jacobian} --- the derivative of the outputs of
    one layer with respect to that layer's inputs. For $\ba = f(\bW\bx)$:
    \[
      \bJ = \pd{\ba}{\bx} = \diag\bigl(f'(\bW\bx)\bigr)\,\bW .
    \]
    This is the quantity connected to the vanishing gradient problem; when
    people say ``the Jacobian'' without further qualification, they usually
    mean this one.
\end{enumerate}
\end{defbox}

\begin{intbox}
The single-layer Jacobian $\diag(f'(\bW\bx))\,\bW$ and the factor
$\bJ\lay{l} = \bW\lay{l}\diag(f'(\bs\lay{l-1}))$ of the backprop recursion
are built from the same two ingredients --- a weight matrix and the
activation derivatives --- just chained in a different order (one goes from
activations to activations, the other from pre-activations to
pre-activations). Either way: backpropagation multiplies by one such matrix
per layer, and its norm decides whether gradients vanish, stay stable or
explode.
\end{intbox}


% ----------------------------------------------------------
\section{Hessian}
\label{sec:hessian}
% ----------------------------------------------------------

The Hessian is the matrix of second derivatives of the loss with respect to
the weights, where all parameters of the network are written as one long
vector $\bw$:
\[
  H_{ij} = \frac{\partial^2 \loss(\by, g(\bx; \bw))}{\partial w_i\,\partial w_j} .
\]
It describes the \emph{curvature} of the loss surface. Its main use is the
second-order Taylor approximation of the error around a point $\bw_0$,
with gradient $\bg = \grad_{\bw}R(\bw_0)$:
\[
  R(\bw) = R(\bw_0) + (\bw - \bw_0)\T\bg
           + \frac12(\bw - \bw_0)\T\bH(\bw - \bw_0)
           + \bigO{\norm{\bw - \bw_0}^3} .
\]
The lecture lists its applications (covered later): \textbf{optimization}
(Newton-type methods --- with a question mark, see below), \textbf{pruning}
(removing the weights whose removal increases the error least),
\textbf{error bars} for predictions, and the \textbf{second-order Taylor
approximation} itself.

\begin{notebox}[title={Why not just use Newton's method for deep networks?}]
For logistic regression, the Hessian is a $D \times D$ matrix and Newton's
method (IRLS) is the standard (\cref{sec:logreg}). A deep network with $W$
weights has a $W \times W$ Hessian: for a million weights that is $10^{12}$
entries, and inverting it costs $\bigO{W^3}$. In addition, the loss of a deep
network is non-convex, so the Hessian is not positive definite in general,
and a Newton step can even go uphill. This is why first-order methods (SGD
and its variants) dominate deep learning.
\end{notebox}


% ----------------------------------------------------------
\section{Efficient Training: Basic Tricks of the Trade}
\label{sec:tricks}
% ----------------------------------------------------------

The last section collects the practical knowledge needed to actually train
an MLP or deep network: how to split the data, how to feed it to the
network, how to prepare the inputs, how to initialize the weights and how to
choose the learning rate and the network size.

\subsection{Training, Validation and Test Set}

\begin{defbox}[title={Training, validation and test set (``sloppy'' definitions)}]
\begin{itemize}
  \item \textbf{Training set} (e.g.\ 80\,\% of the data): the samples that
    are presented to the network; the \emph{weights} are adjusted to fit
    this set.
  \item \textbf{Validation set} (e.g.\ 10\,\%): withheld from the network; it
    is used to check the network during development, and the
    \emph{architecture and hyperparameters} are adjusted on this set.
  \item \textbf{Test set} (e.g.\ 10\,\%): withheld from the network
    \emph{and} from architecture and hyperparameter selection; it is used
    only to check the final trained network.
\end{itemize}
Precise definitions follow in \cref{ch:generalization}.
\end{defbox}

\begin{notebox}[title={Why a separate test set?}]
Every decision made by looking at a data set (stopping epoch, learning
rate, number of layers, \ldots) slightly fits the model to that data set.
After tuning on the validation set, the validation error is therefore an
\emph{optimistic} estimate of the error on new data. Only a test set that
was never used for any decision gives an unbiased estimate.
\end{notebox}

\subsection{Online, Stochastic and Batch Training}

\begin{itemize}
  \item \textbf{Online learning:} a single sample is propagated through the
    network, then the parameters are updated:
    $\bw^{\text{new}} = \bw^{\text{old}} - \eta\,\grad_{\bw}\Remp(y, \bx, \bw)|_{\bw^{\text{old}}}$.
  \item \textbf{Full-batch learning:} all training samples are propagated
    through the network, then one update is performed:
    $\bw^{\text{new}} = \bw^{\text{old}} - \eta\,\grad_{\bw}\Remp(\by, \bX, \bw)|_{\bw^{\text{old}}}$.
  \item \textbf{Stochastic (mini-batch) learning:} a small random subset is
    propagated, then the parameters are updated with the mini-batch estimate
    $\tilde\grad_{\bw}$ of the full-batch gradient:
    \[
      \grad_{\bw}\Remp(\by, \bX, \bw)
      = \frac1N\sum_{n=1}^{N}\grad_{\bw}\loss(y^n, g(\bx^n; \bw))
      \;\approx\;
      \frac1B\sum_{b=1}^{B}\grad_{\bw}\loss(y^{n_b}, g(\bx^{n_b}; \bw))
      = \tilde\grad_{\bw}\Remp .
    \]
\end{itemize}

\begin{intbox}
Full-batch gradients are exact but expensive: one step costs a pass over the
whole data set. Online updates are cheap but very noisy. Mini-batches are
the compromise: the gradient estimate is good enough to make progress,
many updates fit into one pass over the data, and a batch of, say, 32
samples is processed in parallel on a GPU about as fast as a single one.
The noise even helps a little to escape poor regions of the loss surface.
\end{intbox}

\subsection{Epochs, Batch Size and Learning Curves}

Typically, mini-batches are drawn randomly \emph{without replacement} from
the training set. When the training set is used up, a new \textbf{epoch}
starts; one epoch means that every training sample has been propagated
through the network once. With $N$ samples and batch size $B$, an epoch
consists of about $N/B$ updates.

Training is monitored with \textbf{learning curves}:
\begin{itemize}
  \item after each mini-batch: the loss on the training (mini-)batch,
    optionally other metrics such as accuracy;
  \item after each epoch (or after a fixed number of updates): the loss and
    other metrics on the validation set.
\end{itemize}
The lecture's example is a DenseNet trained on CIFAR-100 for 300 epochs:
training and test error drop quickly at first, then flatten out; at epoch 150
both drop sharply once more (at that point the learning rate is decreased,
a common schedule); afterwards the training error goes to almost zero while
the test error stays at about $0.26$. The gap between the two curves is a
first glimpse of \emph{overfitting} (\cref{ch:generalization}).

\subsection{Normalizing Inputs and Outputs}

\begin{itemize}
  \item \textbf{Inputs should be normalized to zero mean and unit variance}:
    compute mean $\mu_d$ and standard deviation $s_d$ of every feature $d$
    \emph{on the training set}, and transform $x_d \mapsto (x_d - \mu_d)/s_d$.
    Training is faster this way. With unnormalized inputs, pre-activations
    can be large, sigmoid/tanh units start in their saturated region where
    the gradient is small; with normalized inputs they start near zero,
    where the gradient is large.
  \item \textbf{Categorical features} are one-hot encoded (\cref{sec:softmax}).
\end{itemize}

\begin{notebox}[title={Use the training statistics everywhere}]
The normalization constants $\mu_d$, $s_d$ are computed on the training set
only and then applied unchanged to validation and test data (and to new
data in deployment). Computing them on the full data set leaks information
from the test set into training.
\end{notebox}

\subsection{Initialization}

\begin{itemize}
  \item Weights should be initialized randomly with \textbf{zero mean}, from
    a Gaussian, uniform or truncated Gaussian distribution.
  \item \textbf{Variance:} $2/J$ for ReLU activations and $1/J$ otherwise,
    where $J$ is the number of inputs of the unit (its \emph{fan-in}, the
    number of units in the layer below).
\end{itemize}
The scaling with $1/J$ keeps the variance of the pre-activations roughly
constant from layer to layer --- which, by the argument of \cref{sec:vgp},
also keeps the delta errors from vanishing or exploding at the start of
training. The factor 2 for ReLU compensates for the fact that a ReLU sets
about half of its inputs to zero. Details in \cref{ch:initialization}.

\begin{notebox}[title={Why not initialize all weights with zero?}]
If all weights of a layer start with the same value, all hidden units of
the layer compute the same function and receive the same delta error ---
they stay identical forever, and the layer behaves like a single unit. The
random initialization is needed to \emph{break the symmetry}.
\end{notebox}

\subsection{Learning Rates}

The learning rate is one of the most important hyperparameters. Typical
values are $0.1$ or $0.01$; it is searched on a logarithmic grid
$10^{-1}, 10^{-2}, 10^{-3}, 10^{-4}, \dots$ A useful heuristic: start
training with a learning rate as high as possible without the network
diverging, and decrease it later (as in the DenseNet learning curve above).

\subsection{Number of Neurons and Hidden Layers}

How large the network has to be depends on the complexity of the problem:
more complex problems need more layers and more neurons. The lecture's
\textbf{rule of thumb}: use as many layers and neurons as needed such that
the network can \emph{overfit} the training data. Then add regularization
(\cref{ch:regularization}) to prevent overfitting.

\begin{figure}[htbp]
\centering
\begin{tikzpicture}
\begin{axis}[axis lines=left, width=0.62\linewidth, height=5cm,
    xmin=0, xmax=10, ymin=0, ymax=1.3, xtick=\empty, ytick=\empty,
    xlabel={model complexity}, ylabel={error},
    label style={font=\small}, samples=120, domain=0.35:9.6,
    axis line style={-{Stealth[length=1.6mm]}}]
  \addplot[very thick, accent] {0.08 + 0.9/(1+1.6*x)};
  \node[font=\scriptsize, accent, anchor=south] at (axis cs:8.6,0.17) {training error};
  \addplot[very thick, figred] {0.08 + 0.9/(1+1.6*x) + 0.012*x^2 - 0.02*x + 0.06};
  \node[font=\scriptsize, figred, anchor=east] at (axis cs:7.6,0.75) {test error};
  \draw[dashed, gray] (axis cs:2.8,0) -- (axis cs:2.8,1.02);
  \node[font=\scriptsize, gray!70!black, align=center, anchor=south] at (axis cs:2.8,1.02) {best complexity};
  \node[font=\scriptsize, gray!70!black, anchor=west] at (axis cs:0.2,0.05) {$\leftarrow$ underfitting};
  \node[font=\scriptsize, gray!70!black, anchor=west] at (axis cs:3.2,0.05) {overfitting $\rightarrow$};
\end{axis}
\end{tikzpicture}
\caption{Outlook on generalization (schematic). The training error keeps
decreasing as the model becomes more complex, but the error on new (test)
data first decreases and then increases again: too complex a model fits the
noise of the training set. Chapter~\ref{ch:generalization} makes this
precise.}
\label{fig:generalization}
\end{figure}


% ----------------------------------------------------------
\section{Excursion: Amari--Hopfield Networks}
\label{sec:hopfield}
% ----------------------------------------------------------

The single-layer lecture closes with a different kind of network.
Amari--Hopfield networks are \textbf{associative memory} systems: they store
a set of patterns and, given a noisy or incomplete version of one of them,
retrieve the complete pattern --- \emph{content-based retrieval}. They are
similar to the Lenz--Ising models of physics, but learnable; Hopfield
received the 2024 Nobel Prize in Physics for them (J.\,J.~Hopfield, 1982),
and he was criticized by J.~Schmidhuber for not crediting S.-I.~Amari, who
had described such networks in 1972.

\begin{defbox}[title={Amari--Hopfield network}]
\begin{itemize}
  \item \textbf{Binary patterns:} inputs and outputs are vectors
    $\bx = (x_1, \dots, x_D) \in \{-1, 1\}^D$; the outputs are fed back as
    new inputs, so the network is \emph{recurrent}.
  \item \textbf{Stored patterns:} $\bx^1, \dots, \bx^N$.
  \item \textbf{Learning (Hebbian rule):}
    \[
      \bW = \frac1N\Bigl(\sum_{n=1}^{N}\bx^n\bx^{n\mathsf{T}}\Bigr) - \bI_D .
    \]
  \item \textbf{Retrieval (inference):} starting from a noisy pattern,
    update
    \[
      \bx^{\text{new}} = \sign\bigl(\bW\bx^{\text{old}} - \bb\bigr).
    \]
  \item \textbf{Energy:} the network is energy-based, with
    \[
      E = -\frac12\bx\T\bW\bx + \bx\T\bb
        = -\frac12\sum_{i=1}^{D}\sum_{j=1}^{D} w_{ij}x_i x_j + \sum_{i=1}^{D} b_i x_i .
    \]
    Stored patterns are (local) minima of the energy; retrieval runs
    ``downhill''.
\end{itemize}
\end{defbox}

\begin{intbox}
The Hebbian rule is Hebb's ``fire together, wire together'' in formula
form: $w_{ij}$ grows when units $i$ and $j$ have the same sign in many stored
patterns ($x^n_i x^n_j = +1$) and shrinks when they have opposite signs.
Since $x_i^2 = 1$, every outer product $\bx^n\bx^{n\mathsf{T}}$ has ones on
the diagonal, so $\frac1N\sum_n\bx^n\bx^{n\mathsf{T}}$ has a diagonal of
exactly one --- subtracting $\bI_D$ removes the self-connections. Each unit
is then pushed towards the sign that the \emph{other} units ``vote'' for,
given the correlations stored in $\bW$.
\end{intbox}

\begin{exbox}[title={Retrieving a pattern with one flipped bit}]
Store two orthogonal patterns in a network with $D = 8$ units,
\[
  \bx^1 = (1, 1, 1, 1, -1, -1, -1, -1), \qquad
  \bx^2 = (1, -1, 1, -1, 1, -1, 1, -1),
\]
with $\bW = \frac12(\bx^1\bx^{1\mathsf{T}} + \bx^2\bx^{2\mathsf{T}}) - \bI_8$
and $\bb = \bm{0}$.
Present $\bx^1$ with its first bit flipped:
$\bx^{\text{old}} = (-1, 1, 1, 1, -1, -1, -1, -1)$.
\begin{itemize}
  \item $\bW\bx^{\text{old}} = (3, 3, 1, 3, -3, -1, -3, -1)$, so
    $\sign(\bW\bx^{\text{old}}) = (1, 1, 1, 1, -1, -1, -1, -1) = \bx^1$: the
    pattern is retrieved in one step.
  \item The energy drops from $E = -6$ for the noisy pattern to $E = -12$ for
    $\bx^1$ (and $\bx^1$, $\bx^2$ are both fixed points of the update).
\end{itemize}
The same happens for every single flipped bit of $\bx^1$.
\end{exbox}


% ----------------------------------------------------------
\section{Chapter Checklist}
% ----------------------------------------------------------

What you should take away from this chapter:
\begin{itemize}
  \item the four single-layer networks (\cref{tab:single-layer}): model,
    loss, and how the loss follows from maximum likelihood (Gaussian noise
    $\to$ squared loss, Bernoulli $\to$ binary cross-entropy, categorical
    $\to$ categorical cross-entropy);
  \item the least-squares solution $\what = (\bX\T\bX)\inv\bX\T\by$, gradient
    descent, and mini-batch SGD;
  \item the perceptron learning rule, its geometric meaning, and the
    convergence theorem with its conditions (linear separability; solution
    not unique; no convergence on non-separable data);
  \item $\sigma' = \sigma(1-\sigma) \le 0.25$; the gradient and Hessian of
    logistic regression; Newton's method and IRLS; why one Newton step
    solves linear regression;
  \item softmax, one-hot encoding, and the common gradient form
    (prediction $-$ target) $\times$ input;
  \item why single-layer networks cannot solve XOR and how a hidden layer
    with a non-linearity does; the universal approximation theorem;
  \item the MLP quantities $s_i$, $a_i$, $f$; common activation functions
    and their derivatives; matching output activations and losses;
  \item backpropagation: the delta errors, the output and hidden-layer
    formulas, $\partial\loss/\partial w_{ij} = \delta_i a_j$, the matrix
    form $\bdelta\lay{l-1} = f'(\bs\lay{l-1}) \odot (\bW\lay{l\mathsf{T}}
    \bdelta\lay{l})$ and $\partial\loss/\partial\bW\lay{l} =
    \bdelta\lay{l}\ba\lay{l-1\mathsf{T}}$; its cost $\bigO{W}$;
  \item why deep networks need non-linearities; the vanishing/exploding
    gradient problem and its remedies;
  \item the three meanings of ``Jacobian'', the Hessian and its uses;
  \item the tricks of the trade: train/validation/test split, online vs.\
    mini-batch vs.\ full-batch, epochs and learning curves, input
    normalization, initialization variance $2/J$ or $1/J$, learning-rate
    heuristics, sizing the network;
  \item Amari--Hopfield networks: Hebbian rule, energy, retrieval.
\end{itemize}


% ==========================================================
%  CHAPTERS 5--10 — OUTLINES
%  Section lists follow the course overview from the first lecture
%  and the forward references on the Chapter 4 slides. Replace
%  \planned{...} blocks with real text once the slides of the
%  corresponding units are available, and adapt the sections to them.
% ==========================================================

% ==========================================================
%  CHAPTER 5 — GENERALIZATION ERROR, TEST SET AND CV
% ==========================================================
\chapter{Generalization Error, Test Set and Cross-Validation}
\label{ch:generalization}

Minimizing the empirical error is not the actual goal --- a network that
memorizes its training data can still be useless on new data. This chapter
makes precise what ``performing well on new data'' means, how to estimate it
honestly, and how to use these estimates to choose models and
hyperparameters.

\planned{Slides not yet available. Expected structure below; adapt it to the
slides.}

\section{Empirical Error and Generalization Error}
\planned{Data-generating distribution; expected loss (generalization error,
risk) $\Rgen$ versus empirical error $\Remp$; why minimizing $\Remp$ is only
a proxy.}

\section{Overfitting and Underfitting}
\planned{Model complexity, capacity; training versus test error
(\cref{fig:generalization}); bias--variance decomposition if covered.}

\section{Estimating the Generalization Error: Test Set}
\planned{Precise definitions of training, validation and test set
(refining \cref{sec:tricks}); unbiasedness of the test-set estimate;
variance of the estimate and test-set size.}

\section{Cross-Validation}
\planned{$k$-fold cross-validation, leave-one-out; using CV for
hyperparameter selection; nested CV; pitfalls (information leakage).}

\section{Chapter Checklist}
\planned{To be written together with the chapter.}


% ==========================================================
%  CHAPTER 6 — CONVOLUTIONAL NEURAL NETWORKS
% ==========================================================
\chapter{Convolutional Neural Networks}
\label{ch:cnn}

Fully connected networks ignore the structure of their input: for them, an
image is just a long vector, and a cat in the top-left corner has nothing to
do with a cat in the bottom-right. Convolutional networks build the spatial
structure into the architecture (\cref{sec:network-types}).

\planned{Slides not yet available. Expected structure below; adapt it to the
slides.}

\section{Convolution and Cross-Correlation}
\planned{Discrete convolution in 1D and 2D; kernels/filters; padding and
stride; output-size formula; channels.}

\section{Convolutional Layers}
\planned{Local connectivity and weight sharing; number of parameters
compared with a fully connected layer; feature maps; translation
equivariance.}

\section{Pooling}
\planned{Max and average pooling; translation invariance; downsampling.}

\section{CNN Architectures and Training}
\planned{Typical layer stacks (convolution $\to$ non-linearity $\to$
pooling, fully connected head); backpropagation through convolution and
pooling; historical architectures (Neocognitron, LeNet, AlexNet) as
covered.}

\section{Chapter Checklist}
\planned{To be written together with the chapter.}


% ==========================================================
%  CHAPTER 7 — LEARNING METHODS: BASIC ALGORITHMS
% ==========================================================
\chapter{Learning Methods: Basic Algorithms}
\label{ch:learning}

\Cref{ch:basic} trained everything with plain (stochastic) gradient descent
and mentioned Newton's method. This chapter looks at the optimization
algorithms actually used to train deep networks.

\planned{Slides not yet available. Expected structure below; adapt it to the
slides.}

\section{Gradient Descent Revisited}
\planned{Batch, mini-batch and online gradient descent
(\cref{sec:tricks}); convergence and the role of the learning rate.}

\section{Momentum and Adaptive Learning Rates}
\planned{Momentum, Nesterov momentum; per-parameter learning rates
(AdaGrad, RMSProp, Adam) as covered.}

\section{Learning-Rate Schedules}
\planned{Step decay (cf.\ the DenseNet learning curve in \cref{sec:tricks}),
warm-up, cosine schedules as covered.}

\section{Second-Order Methods}
\planned{Newton's method (\cref{sec:logreg}), why it does not scale to deep
networks (\cref{sec:hessian}), approximations if covered.}

\section{Chapter Checklist}
\planned{To be written together with the chapter.}


% ==========================================================
%  CHAPTER 8 — REGULARIZATION
% ==========================================================
\chapter{Regularization}
\label{ch:regularization}

The rule of thumb of \cref{sec:tricks} was: make the network large enough to
overfit, then regularize. This chapter is about the second half of that
sentence.

\planned{Slides not yet available. Expected structure below; adapt it to the
slides.}

\section{Weight Decay ($L_2$ and $L_1$ Penalties)}
\planned{Adding a penalty to the empirical error; probabilistic
interpretation as a prior; effect on separable data (\cref{sec:logreg}).}

\section{Early Stopping}
\planned{Monitoring the validation loss (learning curves,
\cref{sec:tricks}); stopping criterion.}

\section{Dropout and Other Noise-Based Methods}
\planned{Dropout at training and test time; data augmentation; other
methods as covered.}

\section{Chapter Checklist}
\planned{To be written together with the chapter.}


% ==========================================================
%  CHAPTER 9 — INITIALIZATION
% ==========================================================
\chapter{Initialization}
\label{ch:initialization}

\Cref{sec:tricks} gave the recipe ``zero mean, variance $2/J$ for ReLU, $1/J$
otherwise''. This chapter derives where these numbers come from.

\planned{Slides not yet available. Expected structure below; adapt it to the
slides.}

\section{Why Initialization Matters}
\planned{Symmetry breaking; connection to vanishing and exploding gradients
(\cref{sec:vgp}).}

\section{Variance-Preserving Initialization}
\planned{Propagation of mean and variance through a layer; LeCun/Glorot
(Xavier) initialization; He initialization for ReLU; forward and backward
view.}

\section{Chapter Checklist}
\planned{To be written together with the chapter.}


% ==========================================================
%  CHAPTER 10 — NORMALIZATION
% ==========================================================
\chapter{Normalization}
\label{ch:normalization}

Normalizing the inputs speeds up training (\cref{sec:tricks}). Normalization
methods apply the same idea to the hidden layers, so that activations stay
well-scaled throughout a deep network.

\planned{Slides not yet available. Expected structure below; adapt it to the
slides.}

\section{Input Normalization Revisited}
\planned{Standardization, whitening; why it helps optimization.}

\section{Batch Normalization and Layer Normalization}
\planned{Normalizing pre-activations with mini-batch or per-sample
statistics; learnable scale and shift; behaviour at training versus test
time.}

\section{Self-Normalizing Networks}
\planned{SELU activation (\cref{tab:activations}) and its fixed point of
zero mean and unit variance; requirements on initialization.}

\section{Chapter Checklist}
\planned{To be written together with the chapter.}


\end{document}
